\documentclass[12pt,a4paper]{article}
\usepackage[utf8]{inputenc}
\usepackage[T1]{fontenc}
\usepackage[ngerman]{babel}
\usepackage{amsmath,amssymb,amsthm}
\usepackage{geometry}
\usepackage{fancyhdr}
\usepackage{titlesec}
\usepackage{enumitem}

\geometry{margin=2.5cm}

\pagestyle{fancy}
\fancyhf{}
\fancyhead[C]{\small Achter Abschnitt \hfill \thepage}
\renewcommand{\headrulewidth}{0.4pt}

\titleformat{\section}{\large\bfseries}{\S.~\thesection.}{0.5em}{}
\titleformat{\subsection}{\normalsize\bfseries}{\thesubsection.}{0.5em}{}

\newcommand{\sect}[1]{\section*{{\S.~#1.}}\addcontentsline{toc}{section}{\S.~#1}}

\newtheorem*{satz}{Satz}

\begin{document}

\title{Lehrbuch der Algebra --- Band~1, Teil~6}
\author{Heinrich Weber}
\date{}
\maketitle

\thispagestyle{empty}
\tableofcontents
\newpage

% §. 59
\sect{59}
\begin{center}
\textbf{Theiler der Ikosaëdergruppe.}
\end{center}
\label{sec:59}
\thispagestyle{fancy}
\fancyhead[C]{\small\S.~59. Theiler der Ikosaëdergruppe. \hfill 229}

Bildet man ausserdem noch $\varphi\,\psi$, $\varphi^{-1}\,\psi$, $\varphi\,\chi\,\psi$, $\chi\,\varphi^{-1}\,\psi$, so folgt:
\[
\varphi\,\psi = \Theta\,\chi\,\psi\,\Theta^{-2},\quad
\varphi^{-1}\psi = \Theta^{-2}\,\chi\,\psi\,\Theta,\quad
\varphi\,\chi\,\psi = \Theta^{-1}\,\chi\,\psi\,\Theta^{2},
\]
\[
\chi\,\varphi^{-1}\,\psi = \Theta^{2}\,\chi\,\psi\,\Theta^{-1},
\]
woraus man sieht, dass man zu keinem anderen Resultate kommen würde, wenn man die Substitution dritter Ordnung, von der man ausgeht, in der zweiten Form $\Theta^{r}\,\chi\,\psi\,\Theta^{s}$ annehmen wollte. Die ganze Gruppe $Q$ ist also durch die angenommene Vierergruppe $D_{2}$ völlig bestimmt, und man erhält sie in der Form
\begin{equation}\label{eq:59.3}
\tag{3}
\begin{aligned}
\varphi &= \Theta\,\chi\,\Theta^{2}, & \varphi\,\psi &= \Theta\,\chi\,\psi\,\Theta^{-2},\\
\varphi^{-1} &= \Theta^{-2}\,\chi\,\Theta^{-1}, & \varphi^{-1}\psi &= \Theta^{-2}\,\chi\,\psi\,\Theta,\\
\varphi\,\chi &= \Theta^{-1}\,\chi\,\Theta^{-2}, & \varphi\,\chi\,\psi &= \Theta^{-1}\,\chi\,\psi\,\Theta^{2},\\
\varphi^{-1}\chi &= \Theta^{2}\,\chi\,\psi\,\Theta^{-1}, & \varphi^{-1}\chi\,\psi &= \Theta^{2}\,\chi\,\Theta.
\end{aligned}
\end{equation}

Man kann diese Gruppe aus den Substitutionen $\varphi$, $\psi$, $\chi$ als den Erzeugenden ableiten und sie in die Form setzen:
\begin{equation}\label{eq:59.4}
\tag{4}
Q = \varphi^{r}\,\chi^{s}\,\psi^{t},\qquad
\begin{aligned}
r &= 0,1,2,\\
s &= 0,1;\quad t = 0,1.
\end{aligned}
\end{equation}

Dass dadurch wirklich eine Tetraëdergruppe dargestellt ist, ergiebt sich aus den Zusammensetzungen, die man mittelst \S.~58, (23), (25) leicht aus (1) findet:
\begin{equation}\label{eq:59.5}
\tag{5}
\begin{aligned}
\psi\,\varphi &= \varphi\,\chi\,\psi, & \chi\,\varphi &= \varphi\,\psi,\\
\psi\,\varphi^{2} &= \chi\,\varphi^{2}\,\psi, &
\chi\,\varphi^{2} &= \varphi^{2}\,\chi\,\psi, &
\psi\,\chi &= \chi\,\psi.
\end{aligned}
\end{equation}

Da in der Tetraëdergruppe fünf Vierergruppen $D_{2}$ enthalten sind, so hat die Ikosaëdergruppe fünf Tetraëdergruppen zu Theilern. Diese können wir aus $Q$ durch Transformation mittelst der Potenzen von $\Theta$ ableiten und erhalten sie in der Form
\[
\Theta^{-r}\,Q\,\Theta^{r}\qquad r = 0,1,2,3,4.
\]

Diese Gruppen haben, ausser der Identität, keine Substitution mit einander gemein. Denn die beiden Gruppen $Q$ und $\Theta^{-1}\,Q\,\Theta$ haben ausser der Identität nur die beiden Elemente
\[
\Theta^{-2}\,\chi\,\Theta^{-1},\quad \Theta\,\chi\,\Theta^{2}
\]
mit einander gemein, und von diesen kommt keines in $\Theta\,Q\,\Theta^{-1}$ vor.

Daraus ergiebt sich nun nach dem Satze 2.~in \S.~28, dass die Ikosaëdergruppe isomorph ist mit einer Permutationsgruppe $60^{\text{sten}}$ Grades von fünf Ziffern. Diese Permutationsgruppe ergiebt sich, wenn wir mit den Nebengruppen
\[
Q,\quad Q\,\Theta,\quad Q\,\Theta^{2},\quad Q\,\Theta^{3},\quad Q\,\Theta^{4}
\]
die sämmtlichen Elemente $\sigma$ der Ikosaëdergruppe verbinden, also
\[
Q\,\sigma,\quad Q\,\Theta\,\sigma,\quad Q\,\Theta^{2}\,\sigma,\quad Q\,\Theta^{3}\,\sigma,\quad Q\,\Theta^{4}\,\sigma
\]
bilden, und die dadurch bewirkte Permutation dieser Nebengruppen untersuchen. Für die erzeugenden Substitutionen der Ikosaëdergruppe $\sigma = \Theta$, $\psi$, $\chi$ erhalten wir so die Permutationen
\[
\begin{pmatrix}
Q, & Q\,\Theta, & Q\,\Theta^{2}, & Q\,\Theta^{3}, & Q\,\Theta^{4}\\
Q, & Q\,\Theta^{4}, & Q\,\Theta^{3}, & Q\,\Theta^{2}, & Q\,\Theta
\end{pmatrix}
\qquad \sigma = \psi
\]
\[
\begin{pmatrix}
Q, & Q\,\Theta, & Q\,\Theta^{2}, & Q\,\Theta^{3}, & Q\,\Theta^{4}\\
Q, & Q\,\Theta^{3}, & Q\,\Theta^{4}, & Q\,\Theta, & Q\,\Theta^{2}
\end{pmatrix}
\qquad \sigma = \chi.
\]
Letzteres findet man aus den Formeln (1) und (2), wonach
\[
\Theta\,\chi = \varphi\,\Theta^{3},\quad
\Theta^{2}\,\chi = \chi\,\varphi^{2}\,\Theta^{4},\quad
\Theta^{3}\,\chi = \varphi^{2}\,\Theta,\quad
\Theta^{4}\,\chi = \varphi\,\chi\,\Theta^{2}.
\]

Man sieht, dass diese Permutationen alle zur ersten Art gehören, und dass also die Permutationsgruppe, um die es sich handelt, keine andere als die alternirende Gruppe von fünf Ziffern (Bd.~I, \S.~177) sein kann. Daraus ergiebt sich auch, dass die Ikosaëdergruppe einfach ist, und folglich mit der Gruppe übereinstimmt, die wir schon in \S.~34 vorläufig als Ikosaëdergruppe bezeichnet haben.

% §. 60
\newpage
\sect{60}
\begin{center}
\textbf{Die Grundformen der Ikosaëdergruppe.}
\end{center}
\label{sec:60}
\fancyhead[C]{\small\S.~60. Grundformen des Ikosaëders. \hfill 231}

Wir haben oben die eine der drei Grundformen der Ikosaëdergruppe $f$ gefunden:
\begin{equation}\label{eq:60.1}
\tag{1}
f = x_{1}\,x_{2}\,(x_{1}^{10} + 11\,x_{1}^{5}\,x_{2}^{5} - x_{2}^{10}).
\end{equation}

Es fehlen uns also noch zwei dieser Formen, eine vom $20^{\text{ten}}$ und eine vom $30^{\text{ten}}$ Grade.

Die Grundform $20^{\text{ten}}$ Grades ergiebt sich als die Hesse'sche Covariante von $f$. Wir setzen sie
\begin{equation}\label{eq:60.2}
\tag{2}
H = \frac{1}{121}\,\bigl[f''(x_{1},x_{1})\,f''(x_{2},x_{2}) - f''(x_{1},x_{2})^{2}\bigr],
\end{equation}
und finden durch einfache Rechnung:
\begin{equation}\label{eq:60.3}
\tag{3}
H = -x_{1}^{20} + x_{2}^{20} + 228\,(x_{1}^{15}\,x_{2}^{5} - x_{1}^{5}\,x_{2}^{15}) - 494\,x_{1}^{10}\,x_{2}^{10}.
\end{equation}

Die Grundform $30^{\text{ten}}$ Grades können wir als die Functional-determinante von $H$ und $f$ definiren. Setzen wir
\begin{equation}\label{eq:60.4}
\tag{4}
T = \frac{\partial f}{\partial x_{1}}\,\frac{\partial H}{\partial x_{2}} - \frac{\partial f}{\partial x_{2}}\,\frac{\partial H}{\partial x_{1}},
\end{equation}
so findet sich
\begin{multline}\label{eq:60.5}
\tag{5}
T = (x_{1}^{30} + x_{2}^{30}) + 522\,(x_{1}^{25}\,x_{2}^{5} - x_{1}^{5}\,x_{2}^{25})\\
- 10005\,(x_{1}^{20}\,x_{2}^{10} + x_{1}^{10}\,x_{2}^{20}).
\end{multline}

Auch zwischen diesen drei Formen besteht eine identische Relation. Um sie zu finden, setzen wir
\[
A = x_{1}^{20} - x_{2}^{20},\quad u = x_{1}^{5}\,x_{2}^{5},
\]
und erhalten, wenn wir bei $T$ den Factor $x_{1}^{2} + x_{2}^{2}$ herausnehmen und dann $T^{2}$ bilden:
\[
\begin{aligned}
f &= x_{1}\,x_{2}\,(A + 11\,u),\\
H &= -A^{2} + 228\,A\,u - 496\,u^{2},\\
T^{2} &= (A^{2} + u^{2})\,(A^{2} + 522\,A\,u - 10006\,u^{2})^{2}.
\end{aligned}
\]
Daraus erhält man durch die numerische Berechnung
\begin{equation}\label{eq:60.6}
\tag{6}
T^{2} + H^{3} = 1728\,f^{5},
\end{equation}
was die gesuchte Relation ist.

Da die Ikosaëdergruppe, wie wir gesehen haben, einfach ist, so kann sie nach \S.~40 keine relativen Invarianten haben. Daraus ergiebt sich, dass die Ikosaëderformen $f$, $H$, $T$ absolut ungeändert bleiben, wenn man sie den mit der Determinante~1 dargestellten binären Ikosaëdersubstitutionen unterwirft. Dasselbe lässt sich aber auch, ohne jene allgemeinen Sätze zu benutzen, in folgender Weise direct nachweisen.

Die Relation (6) lässt sich in der Form darstellen:
\begin{equation}\label{eq:60.7}
\tag{7}
\frac{T^{2}}{f^{5}} + \frac{H^{3}}{f^{5}} = 1728,
\end{equation}
und die beiden Quotienten $T^{2}\!:f^{5}$, $H^{3}\!:f^{5}$ können wir als Functionen der einen Veränderlichen $z = x_{1}\!:x_{2}$ auffassen. Wenden wir auf diese Variable eine Substitution der Ikosaëdergruppe an, so ändern sich diese beiden Quotienten nur um constante Factoren, d.~h. wenn wir
\begin{equation}\label{eq:60.8}
\tag{8}
\Phi(x) = \frac{H^{3}}{f^{5}},\quad \Psi(x) = \frac{T^{2}}{f^{5}}
\end{equation}
setzen und mit $h$, $k$ zwei Constanten bezeichnen, so ist
\[
\Phi(y) = h\,\Phi(x),\quad \Psi(y) = k\,\Psi(x),
\]
vorausgesetzt dass $\binom{x_{1}}{x_{2}}$ eine der Ikosaëdersubstitutionen ist. Da nun aber $y$ sowohl als $x$ eine unabhängige Variable ist, so besteht nach (7) die Identität
\[
\Phi(x) + \Psi(x) = h\,\Phi(x) + k\,\Psi(x) = 1728,
\]
und da $\Phi$ und $\Psi$ nicht in constantem Verhältnisse stehen, so muss $h = k = 1$ sein, d.~h. die Quotienten $\Phi$ und $\Psi$ bleiben bei den Ikosaëdersubstitutionen absolut ungeändert.

Daraus ergiebt sich weiter, dass $f$, $H$, $T$ bei den homogenen Ikosaëdersubstitutionen mit der Determinante~1 absolut ungeändert bleiben.

Denn wenn $f$ durch eine solche Substitution in $c\,f$ übergeht, so gehen $H$ und $T$ in $c^{2}H$ und $c^{3}T$ über [nach (2) und (4)], und die Quotienten $\Phi$ und $\Psi$ in $c^{0}$ und $c^{0}$. Da andererseits $\Phi$ und $\Psi$ ungeändert bleiben, so ist $c = 1$.

Setzen wir $\Psi(x) = z$, so wird $\Phi(x) = 1728 - z$, und wir erhalten aus (8) die beiden Gleichungen:
\begin{align}
\label{eq:60.9}\tag{9} H^{3} - z\,f^{5} &= 0,\\
\label{eq:60.10}\tag{10} T^{2} - (1728 - z)\,f^{5} &= 0,
\end{align}
von denen wegen der Identität (6) die eine aus der anderen folgt, so dass es eigentlich nur zwei verschiedene Formen für eine und dieselbe Gleichung sind. Betrachten wir nun darin $z$ als gegeben, so haben wir eine Gleichung $60^{\text{ten}}$ Grades für $z$, die die \emph{Ikosaëdergleichung} heisst.

Auf die Eigenschaften dieser Gleichung und ihre Beziehung zu der allgemeinen Gleichung $5^{\text{ten}}$ Grades kommen wir in einem späteren Abschnitte zurück.



% §. 61
\newpage
\sect{61}
\begin{center}
\textbf{Die Invarianten des Ikosaëders.}
\end{center}
\label{sec:61}
\fancyhead[C]{\small\S.~61. Invarianten des Ikosaëders. \hfill 233}

Durch die Grundformen der Polyëdergruppen, die wir bisher kennen gelernt haben, ist, wie wir jetzt beweisen können, das Gebiet der Invarianten der betreffenden Gruppen erschöpft. Wir beschränken uns bei diesem Beweise auf die Betrachtung der Ikosaëdergruppe, da für die anderen Gruppen ganz ähnliche Schlüsse zu machen sind, die wir dem Leser überlassen können. Der Umstand, dass bei der Tetraëder- und Octaëdergruppe neben den absoluten auch relative Invarianten vorkommen, während die Ikosaëdergruppe als einfache Gruppe nur absolute Invarianten hat, ist hierbei von keinem wesentlichen Einflusse.

Wir beweisen also, dass alle Invarianten der Ikosaëdergruppe sich als ganze rationale Functionen der drei Formen $f$, $H$, $T$ darstellen lassen.

Dazu führt uns folgende Schlusskette:

\medskip\noindent
\textbf{1.} \emph{Keine, zwei der Ikosaëderformen $f$, $H$, $T$ haben einen gemeinschaftlichen Theiler.}

Dies folgt unmittelbar aus der Definition dieser Functionen, wonach die Gleichungen $f = 0$, $H = 0$, $T = 0$ die fünfzähligen, dreizähligen und zweizähligen Pole liefern, die alle von einander verschieden sind.

\medskip\noindent
\textbf{2.} \emph{Eine Doppelwurzel der Ikosaëdergleichung}
\begin{equation}\label{eq:61.1}
\tag{1}
H^{3} - z\,f^{5} = 0
\end{equation}
\emph{ist nothwendig eine Wurzel von $f$, $H$ oder $T$, und kann also nur für einen der Werthe $z = 0$, $\infty$, $1728$ eintreten.}

Wenn nämlich (1) eine Doppelwurzel hat, so müssen mit der Function zugleich die erste Ableitung, oder, wenn man die homogene Form anwendet, nach der Euler'schen Theorie [Bd.~I, \S.~17,~(5)] die beiden Ableitungen nach $x_{1}$ und $x_{2}$ zugleich verschwinden, also
\[
3H^{2}\,H'(x_{1}) - 5z\,f^{4}\,f'(x_{1}) = 0,\qquad
3H^{2}\,H'(x_{2}) - 5z\,f^{4}\,f'(x_{2}) = 0;
\]
folglich bleiben, da $H$ und $f$ nach 1.~nicht zugleich verschwinden, nur drei Möglichkeiten: entweder $H = 0$, $z = 0$, oder $f = 0$, $z = \infty$, oder endlich
\[
H'(x_{1})\,f'(x_{2}) - H'(x_{2})\,f'(x_{1}) = 0,
\]
d.~h. $T = 0$, $z = 1728$ [\S.~60,~(4)]. Hieran schliesst sich auch der evidente Satz:

\medskip\noindent
\textbf{3.} \emph{Wenn $J(x_{1},x_{2})$ irgend eine invariante Form der Ikosaëdergruppe ist, und $\xi$ eine ihrer Wurzeln, so dass $J(\xi,1) = 0$ ist, so sind auch alle die Grössen Wurzeln von $J$, die aus $\xi$ durch die gebrochenen Ikosaëdersubstitutionen hervorgehen.}

Wenn daher $J$ mit einer der Functionen $f$, $H$, $T$ einen Theiler gemein hat, so ist $J$ durch die betreffende Form theilbar.

Wir denken uns nun zunächst aus $J$ möglichst hohe Potenzen der drei Grundformen $f$, $H$, $T$ weggehoben, so dass $J$ zu diesen Functionen theilerfremd ist. Ist dann $\xi$ eine Wurzel von $J$, so können wir $\eta$ so bestimmen, dass $\xi$ eine Wurzel der Form
\[
\vartheta = H - \eta\,f
\]
ist, und $\vartheta$ ist dann nach 2.~eine einfache Wurzel von $\vartheta$. Es müssen dann auch nach 3.~alle übrigen Wurzeln von $\vartheta$ zugleich Wurzeln von $J$ sein, d.~h. $J$ ist durch $\vartheta$ theilbar. Der Quotient ist wieder eine invariante Form des Ikosaëders, und der Schluss lässt sich wiederholen. Wir kommen so zu dem Satze:

\medskip\noindent
\textbf{4.} \emph{Jede Invariante des Ikosaëders lässt sich in der Form darstellen:}
\begin{equation}\label{eq:61.2}
\tag{2}
J(x_{1},x_{2}) = C\,f^{\alpha}\,H^{\beta}\,T^{\gamma}\,F(\Phi,\Psi),
\end{equation}
\emph{worin $\alpha$, $\beta$, $\gamma$ nicht negative ganze Zahlen, $C$ eine Constante und $F$ eine ganze homogene Function bedeuten.}

% §. 62
\newpage
\sect{62}
\begin{center}
\textbf{Polyëdergruppen der zweiten Art.\\Krystallographische Gruppen.}
\end{center}
\label{sec:62}
\fancyhead[C]{\small\S.~62. Gruppen der zweiten Art. \hfill 235}

Wir haben schon im \S.~51 auf Gruppen linearer ternärer Substitutionen aufmerksam gemacht, die ausser den eigentlichen auch uneigentlich orthogonale Substitutionen enthalten, und die wir \emph{Gruppen der zweiten Art} nennen.

Wir wollen die endlichen unter ihnen jetzt als \emph{Polyëdergruppen der zweiten Art} oder auch als \emph{erweiterte Polyëdergruppen}, und die darin enthaltenen Substitutionen mit der Determinante $-1$ als \emph{Substitutionen der zweiten Art} bezeichnen.

Nach den Resultaten des \S.~50 sind diese Gruppen isomorph mit den Gruppen binärer linearer Substitutionen mit der Determinante $+1$ und $-1$, während bei den gebrochenen Substitutionen, die uns auf die Polyëdergruppen geführt haben, die beiden Arten nicht unterscheidbar sind.

Die endlichen Gruppen der zweiten Art sind, abgesehen von ihrem allgemeinen gruppentheoretischen Interesse, wichtig wegen ihrer Anwendung in der Krystallographie. Ihre vollständige Bestimmung bietet jetzt keine wesentlichen Schwierigkeiten mehr.

Wir verstehen unter Substitutionen schlechweg binäre lineare Substitutionen mit der Determinante $\pm 1$, und erinnern daran, dass zwei solche Substitutionen, deren Coefficienten sich nur durch das gemeinschaftliche Vorzeichen unterscheiden, wie
\[
\begin{pmatrix}\alpha&\beta\\\gamma&\delta\end{pmatrix},
\quad
\begin{pmatrix}-\alpha&-\beta\\-\gamma&-\delta\end{pmatrix}
\]
als nicht verschieden gelten (\S.~49).

Da eine Gruppe linearer Substitutionen die identische Substitution enthält, die von der Determinante $+1$ ist, so müssen in jeder Gruppe $Q$ der zweiten Art neben den uneigentlichen auch eigentliche Substitutionen vorkommen, und diese eigentlichen Substitutionen bilden für sich eine Gruppe $P$. Ist $\varrho$ irgend eine in $Q$ vorkommende Substitution der zweiten Art, so kommt jede Composition von $\varrho$ mit einer Substitution aus $\varrho$ der zweiten Art in $P$ vor, und daher können wir $Q$ immer so zerlegen:
\begin{equation}\label{eq:62.1}
\tag{1}
Q = P + P\,\varrho.
\end{equation}

Die Gruppe $P$ muss eine der früher betrachteten Polyëdergruppen sein. Es muss $\varrho^{-1}\,P\,\varrho = P$ sein, und $P$ ist also ein Normaltheiler von $Q$. Ist umgekehrt $P$ eine Polyëdergruppe und $\varrho$ eine Substitution zweiter Art, die der Bedingung $\varrho^{-1}\,P\,\varrho = P$ genügt, so ist $Q = P + P\,\varrho$ eine Gruppe der zweiten Art.

Auch hier betrachten wir zwei Gruppen, die durch Transformation aus einander hervorgehen, als nicht wesentlich verschieden. Wir haben, um alle $Q$ zu bilden, die verschiedenen Fälle durchzugehen.

\medskip\noindent
\textbf{I.} Ist $P$ die Einheitsgruppe, die wir als cyklische Gruppe ersten Grades mit $C_{1}$ bezeichnen, besteht also $P$ nur aus der identischen Substitution, so muss $\varrho^{2} = 1$ sein, und dies ist auch ausreichend. Wir können hier durch Transformation $\varrho$ auf die Form $\binom{0\;\alpha}{\beta\;0}$ bringen und erhalten als Bedingung $\alpha\beta = \pm 1$, und demnach zwei Formen der Substitution $\varrho$:
\[
\varrho' = \begin{pmatrix}0&1\\1&0\end{pmatrix},\qquad
\varrho'' = \begin{pmatrix}0&1\\-1&0\end{pmatrix},
\]
von denen der erste die \emph{Inversion}, der zweite die \emph{Spiegelung} genannt wird. Ueberträgt man sie nach \S~50,~(1) auf eine rechtwinkelige Coordinatentransformation, so bedeutet $\varrho'$ die gleichzeitige Vorzeichenänderung aller drei Coordinaten, $\varrho''$ die Vertauschung von $x$ mit $-z$. Es ergeben sich hieraus die beiden Gruppen der zweiten Art:
\begin{equation}\label{eq:62.2}
\tag{2}
C_{1}' = \begin{pmatrix}1&0\\0&1\end{pmatrix},\;
\begin{pmatrix}1&0\\0&-1\end{pmatrix};\qquad
C_{1}'' = \begin{pmatrix}1&0\\0&1\end{pmatrix},\;
\begin{pmatrix}-1&0\\0&-1\end{pmatrix},\qquad h = 0,1.
\end{equation}

\medskip\noindent
\textbf{II.} Es sei $P$ eine cyklische Gruppe $C_{n}$, $n > 1$.

Setzen wir
\begin{equation}\label{eq:62.3}
\tag{3}
c = \begin{pmatrix}\varepsilon&0\\0&\varepsilon^{-1}\end{pmatrix},\qquad
\varepsilon = e^{\frac{2\pi i}{n}},
\end{equation}
so ist nach \S.~55:
\begin{equation}\label{eq:62.4}
\tag{4}
C_{n} = 1,\; c,\; c^{2},\; \ldots,\; c^{n-1},
\end{equation}
worin
\[
c = \begin{pmatrix}\varepsilon&0\\0&\varepsilon^{-1}\end{pmatrix}.
\]
Wenn nun
\[
\varrho = \begin{pmatrix}\alpha&\beta\\\gamma&\delta\end{pmatrix}
\]
ist, so bilden wir zunächst
\begin{align}
\label{eq:62.5}\tag{5}
\varrho^{-1}\,c\,\varrho &=
\begin{pmatrix}
\alpha\delta\,\varepsilon - \beta\gamma\,\varepsilon^{-1},&
\alpha\beta\,(\varepsilon - \varepsilon^{-1})\\
-\gamma\delta\,(\varepsilon - \varepsilon^{-1}),&
-\beta\gamma\,\varepsilon + \alpha\delta\,\varepsilon^{-1}
\end{pmatrix},\\
\label{eq:62.6}\tag{6}
\varrho^{2} &=
\begin{pmatrix}
\alpha^{2} + \beta\gamma,&
(\alpha + \delta)\,\beta\\
(\alpha + \delta)\,\gamma,&
\gamma\beta + \delta^{2}
\end{pmatrix}.
\end{align}

Diese beiden Substitutionen müssen in der Reihe (4) vorkommen, und es muss also $\beta\gamma = 0$, $\alpha\delta = 0$ sein; also müssen entweder $\beta$ und $\gamma$ oder $\alpha$ und $\delta = 0$ sein. Ist $\beta = \gamma = 0$, so folgt aus (6), dass $\alpha$ von der Form $\varepsilon^{r\pm\frac{n}{2}}$, also
\[
\varrho = \begin{pmatrix}\varepsilon^{r\pm\frac{n}{2}}&0\\0&\varepsilon^{-(r\pm\frac{n}{2})}\end{pmatrix}
\]
sein muss, worin $r$ eine ganze Zahl bedeutet. Je nachdem $r$ ungerade oder gerade angenommen wird, erhalten wir zwei verschiedene Gruppen, die sich nicht durch Transformation auf einander zurückführen lassen:
\[
\text{1)}\quad
\begin{pmatrix}\varepsilon^{\frac{h}{2}}&0\\0&(-1)^{h}\,\varepsilon^{-\frac{h}{2}}\end{pmatrix},\quad
h = 0,1,\ldots,2n-1;
\]
\[
\text{2)}\quad
\begin{pmatrix}\varepsilon^{h}&0\\0&\pm\varepsilon^{-h}\end{pmatrix},\quad
h = 0,1,\ldots,n-1.\quad \text{(beide Vorzeichen)}
\]
Wenn $n$ gerade ist, so kommen $C_{n}'$ und $C_{n}''$ beide unter $C_{\frac{n}{2}}$ vor; ist aber $n$ ungerade, so kommt $C_{n}'$ in $C_{2n}$, $C_{n}''$ in $C_{n}$ vor.

Ist sodann $\alpha = \delta = 0$, $\beta\gamma = 1$, so können wir $\varrho$ durch die Transformation
\[
\begin{pmatrix}\delta&0\\0&1\end{pmatrix}\,
\begin{pmatrix}1&0\\0&\delta\end{pmatrix}\,
\begin{pmatrix}\alpha&\beta\\\gamma&\delta\end{pmatrix}\,
\begin{pmatrix}\delta^{-1}&0\\0&1\end{pmatrix}\,
\begin{pmatrix}1&0\\0&\delta\end{pmatrix}
= \begin{pmatrix}0&\beta\\\gamma&0\end{pmatrix},
\]
durch die die Gruppe $P$ ungeändert bleibt, umformen, und erhalten noch eine dritte Gruppe:
\[
\text{3)}\quad
C_{n}''' =
\begin{pmatrix}\varepsilon^{h}&0\\0&\varepsilon^{-h}\end{pmatrix},\;
\begin{pmatrix}0&\varepsilon^{h}\\\varepsilon^{-h}&0\end{pmatrix},\qquad
h = 0,1,\ldots,n-1.
\]

\medskip\noindent
\textbf{III.} Es sei $P$ die Diëdergruppe $D_{n}$, die aus den Substitutionen
\begin{equation}\label{eq:62.8}
\tag{8}
\begin{aligned}
D_{n} = 1,\; c,\; c^{2},\; \ldots,\; c^{n-1}\\
d,\; cd,\; c^{2}d,\; \ldots,\; c^{n-1}d = \begin{pmatrix}0&\varepsilon^{h}\\\varepsilon^{-h}&0\end{pmatrix}
\end{aligned}
\end{equation}
besteht. Die Substitution $c$ ist vom $n^{\text{ten}}$, die $c^{h}d$ sind alle vom $2^{\text{ten}}$ Grade. Es ist aber $\varrho^{-1}\,c\,\varrho$ vom $n^{\text{ten}}$ Grade, und muss daher, wenn $n > 2$ ist, unter den Potenzen von $c$ enthalten sein. Folglich sind zur Bestimmung von $\varrho$ die vorigen Betrachtungen anwendbar, und in der Gruppe $Q$ muss eine der Gruppen $C_{n}'$, $C_{n}''$, $C_{n}'''$ enthalten sein. Da $d$ nicht in diesen Gruppen vorkommt, so ergeben sich für die Diëdergruppen zweiter Art die Formen
\begin{equation}\label{eq:62.9}
\tag{9}
C_{n}' + C_{n}'d,\quad C_{n}'' + C_{n}''d,\quad C_{n}''' + C_{n}'''d,
\end{equation}
von denen die beiden ersten
\[
\text{1)}\quad
D_{n}' =
\begin{pmatrix}\varepsilon^{\frac{h}{2}}&0\\0&(-1)^{h}\,\varepsilon^{-\frac{h}{2}}\end{pmatrix},\;
\begin{pmatrix}0&\varepsilon^{\frac{h}{2}}\\(-1)^{h}\,\varepsilon^{-\frac{h}{2}}&0\end{pmatrix},\quad
h = 0,1,\ldots,2n-1;
\]
\[
\text{2)}\quad
D_{n}'' =
\begin{pmatrix}\varepsilon^{h}&0\\0&\pm\varepsilon^{-h}\end{pmatrix},\;
\begin{pmatrix}0&\varepsilon^{h}\\\pm\varepsilon^{-h}&0\end{pmatrix},\qquad
h = 0,1,\ldots,n-1
\]
sind, während die dritte Gruppe (9)
\[
\begin{pmatrix}\varepsilon^{h}&0\\0&\varepsilon^{-h}\end{pmatrix},\;
\begin{pmatrix}0&\varepsilon^{h}\\\varepsilon^{-h}&0\end{pmatrix}
\]
\[
\begin{pmatrix}\varepsilon^{h}&0\\0&-\varepsilon^{-h}\end{pmatrix},\;
\begin{pmatrix}0&\varepsilon^{h}\\-\varepsilon^{-h}&0\end{pmatrix}
\]
\[
h = 0,1,\ldots,n-1
\]
ergibt, was bei geradem $n$ mit $D_{\frac{n}{2}}'$, bei ungeradem $n$ mit $D_{n}''$ übereinstimmt.

Der Fall $n = 2$, wo $c = \binom{0\;1}{-1\;0}$ ist, bildet nur eine scheinbare Ausnahme, weil in diesem Falle
\begin{equation}\label{eq:62.10}
\tag{10}
\varrho^{-1}\,c\,\varrho = \pm d\quad\text{oder}\quad
\varrho^{-1}\,c\,\varrho = \mp cd
\end{equation}
sein kann. Verfolgt man diese Annahmen durch einfache Rechnung, so ergeben sich noch zwei Gruppen:
\begin{equation}\label{eq:62.11}
\tag{11}
\begin{aligned}
&1,\; c,\; d,\; cd,\; \varrho_{1},\; \varrho_{1}c,\; \varrho_{1}d,\; \varrho_{1}cd,\\
&1,\; c,\; d,\; cd,\; \varrho_{2},\; \varrho_{2}c,\; \varrho_{2}d,\; \varrho_{2}cd,
\end{aligned}
\end{equation}
worin
\[
\varrho_{1} = \begin{pmatrix}\frac{1}{\sqrt{2}}&\frac{1}{\sqrt{2}}\\[-2pt]\frac{1}{\sqrt{2}}&-\frac{1}{\sqrt{2}}\end{pmatrix},\qquad
\varrho_{2} = \begin{pmatrix}\frac{i}{\sqrt{2}}&\frac{1}{\sqrt{2}}\\[-2pt]\frac{1}{\sqrt{2}}&-\frac{i}{\sqrt{2}}\end{pmatrix}.
\]

\newpage
\fancyhead[C]{\small\S.~62. Krystallographische Gruppen. \hfill 239}

Denn aus den Formeln
\[
\varrho_{1}\,c = d\,\varrho_{1},\quad
\varrho_{1}\,d = c\,\varrho_{1},\quad
\varrho_{1}\,cd = -cd\,\varrho_{1},\quad
\varrho_{2}\,c = -d\,\varrho_{2},\quad
\varrho_{2}\,d = c\,\varrho_{2},\quad
\varrho_{2}\,cd = cd\,\varrho_{2}
\]
und aus diesen Formeln folgt, dass die Gruppen (11) durch Transformation mit $\varrho_{1}$ und $\varrho_{2}$ in $D_{2}$ transformirt werden.

Bei der Bestimmung der übrigen Polyëdergruppen zweiter Art sind die folgenden allgemeinen Bemerkungen von Nutzen.

Die Inversion $j = \binom{1\;\phantom{-}0}{0\;-1}$ ist eine Aehnlichkeitssubstitution (\S.~37).

Es ist aber
\[
j^{-1}\,\begin{pmatrix}\alpha&\beta\\\gamma&\delta\end{pmatrix}\,j
= \begin{pmatrix}\alpha&-\beta\\-\gamma&\delta\end{pmatrix},
\]
und daher mit jeder anderen Substitution vertauschbar, und folglich können wir aus jeder Polyëdergruppe erster Art wenigstens eine Gruppe zweiter Art herleiten:
\begin{equation}\label{eq:62.12}
\tag{12}
P + P\,j.
\end{equation}

Um die anderen etwa noch vorhandenen Gruppen dieser Art zu finden, erinnern wir uns, dass nach dem Sylow'schen Satze (\S.~29,~I.) in jeder Gruppe $G$ eine Gruppe enthalten sein muss, deren Grad die höchste Potenz von~2 ist, die im Grade von~$G$ aufgeht. Die erweiterten Tetraëder-, Octaëder- und Ikosaëdergruppen haben aber die Grade 24, 48, 120, müssen also einen Theiler vom Grade 8, 16, 8 enthalten.

Es sei nun $T$ die Tetraëder-, $O$ die Octaëder- und $J$ die Ikosaëdergruppe. In $T$ und $J$ ist eine Vierergruppe $D_{2}$ enthalten, aber keine Substitution vierter Ordnung, in $O$ eine Diëdergruppe $D_{4}$ und keine Substitution achter Ordnung, und es muss daher unter den erweiterten Polyëdergruppen eine der unter III. betrachteten erweiterten Diëdergruppen enthalten sein.

Von den erweiterten Diëdergruppen entsteht aber $D_{2}'$ aus $D_{2}$ durch Inversion, während
\[
D_{2}'' = D_{2} + D_{2}\,\varrho''
\]
durch Composition von $D_{2}$ mit
\[
\varrho'' = \begin{pmatrix}0&-i\\i&\phantom{-}0\end{pmatrix}
\]
entsteht, also
\[
D_{2}'' = D_{2} + D_{2}\,j.
\]
Nun ist allgemein:
\[
\varrho''^{-1}\,\begin{pmatrix}\alpha&\beta\\\gamma&\delta\end{pmatrix}\,\varrho''
= \begin{pmatrix}\delta&-\gamma\\-\beta&\alpha\end{pmatrix};
\]
und daraus schliesst man nach \S.~56,~(4), dass $\varrho''^{-1}\,T\,\varrho'' = T$ ist, dass also aus der Tetraëdergruppe zwei erweiterte Gruppen $T' = T + T\,j$ und $T'' = T + T\,\varrho''$ abgeleitet werden können; aber auch nur zwei, weil nach \S.~56 die Tetraëdergruppe durch eine darin enthaltene Vierergruppe völlig bestimmt ist. Denn wenn wir statt der Erweiterung $D_{2}'$ die Erweiterung der Vierergruppe zu einer der Gruppen (11) wählen, so erhalten wir eine erweiterte Tetraëdergruppe $T'''$, die durch Transformation mit $\varrho_{1}$ oder $\varrho_{2}$ in $T'$ übergeht.

Da die Substitution
\[
\varrho''' = \begin{pmatrix}0&1\\1&0\end{pmatrix}
\]
mit der Octaëdergruppe \S.~57,~(8) nicht vertauschbar ist, so lässt sich aus der Octaëdergruppe ausser durch Inversion keine weitere Gruppe zweiter Art ableiten, und dasselbe gilt von der Ikosaëdergruppe. Man muss, um diese Schlussweise anzuwenden, die Ikosaëdergruppe so transformiren, dass eine der darin enthaltenen Vierergruppen die einfachste Gestalt annimmt, was etwas weitläufig, aber durchaus nicht schwierig ist, und hier nicht weiter ausgeführt werden soll.

Wir erhalten also noch folgende Gruppen zweiter Art:
\[
\begin{aligned}
\text{IV.}\quad &\text{1)}\; T' = T + T\,j,\qquad \text{2)}\; T'' = T + T\,\varrho'',\\
\text{V.}\quad &\; O' = O + O\,j,\\
\text{VI.}\quad &\; J' = J + J\,j.
\end{aligned}
\]

\newpage
\fancyhead[C]{\small\S.~62. Krystallographische Gruppen. \hfill 241}

Hierin sind die 32 Symmetriesysteme der Krystallographen enthalten. Es sind die Gruppen:
\[
\begin{array}{lll}
C_{1}' & C_{1}'' & 1,\\
C_{2}' & C_{2}'' & C_{2}''',\\
C_{3}' & C_{3}'' & C_{3}''',\\
C_{4}' & C_{4}'' & C_{4}''',\\
C_{6}' & C_{6}'' & C_{6}''',\\
D_{2}' & D_{2}'' & D_{2}'''\\
D_{3}' & D_{3}'' & D_{3}'''\\
D_{4}' & D_{4}'' & D_{4}'''\\
D_{6}' & D_{6}'' & D_{6}'''\\
T,\; T',\; T'' & O,\; O' & J'
\end{array}
\]

Die übrigen Polyëdergruppen sind in der Krystallographie durch das krystallographische Gesetz der rationalen Indices ausgeschlossen\footnote{Vgl.~Schönflies, Krystallsysteme u.~Krystallstructur. Leipzig 1891.}.

\begin{flushright}
Weber, Algebra.~II.
\end{flushright}



% Neunter Abschnitt
\newpage
\begin{center}
\Large\textbf{Neunter Abschnitt.}\\[4pt]
\Large\textbf{Congruenzgruppen.}
\end{center}
\label{sec:neunter}
\thispagestyle{fancy}
\fancyhead[C]{\small Neunter Abschnitt. \hfill \thepage}

% §. 63
\sect{63}
\begin{center}
\textbf{Functionen-Congruenzen.}
\end{center}
\label{sec:63}
\fancyhead[C]{\small\S.~63. Functionen-Congruenzen. \hfill 243}

Aus den linearen Substitutionen, deren Bildung und Zusammensetzung wir in den früheren Abschnitten kennen gelernt haben, lassen sich noch eine ganz andere Art endlicher Gruppen ableiten, die \emph{Congruenzgruppen}.

Diese Congruenzgruppen haben ein mannigfaches Interesse. Sie stammen zunächst aus der Theorie der elliptischen Functionen, wo sie sich als Galois'sche Gruppen der Modulargleichungen einstellen. Sie sind aber auch für die allgemeine Gruppentheorie von Wichtigkeit, weil sie uns ein Mittel geben, ganze Reihen von einfachen Gruppen zu bilden. Dies ist um so bemerkenswerther, als, wie wir im vierten Abschnitte gesehen haben, die einfachen Gruppen, wenigstens unter den niedrigeren Gradzahlen, sehr selten sind.

Die Theorie dieser Congruenzgruppen wird nicht nur ausserordentlich verallgemeinert, sondern die herrschenden Gesetze treten weit schärfer hervor, wenn man sich auf eine von Gauss herrührende Erweiterung des Congruenzbegriffes stützt, die seitdem mehrfach für die Probleme der Gruppentheorie, besonders von Galois, ausgebildet und angewandt worden ist\footnote{Galois, \,Sur la théorie des nombres\,. Bulletin des sciences math.\ de Ferussac. 1830. (Vergl.\ die Note zu \S.~151 des ersten Bandes.) --- Schönemann, Grundzüge einer allgemeinen Theorie der höheren Congruenzen, deren Modulus eine reelle Primzahl ist. (Crelle's Journ.\ f.\ Mathematik, Bd.~31, 1846.) --- Dedekind, Abriss einer Theorie der höheren Congruenzen in Bezug auf einen reellen Primzahlmodulus. (Crelle's Journ.\ f.\ Mathematik, Bd.~54, 1856.)}.

\medskip\noindent\textbf{Um eine Grundlage für diese Theorie zu gewinnen}, betrachten wir eine ganze Function einer veränderlichen Grösse $t$
\begin{equation}\label{eq:63.1}
\tag{1}
f(t) = a_{0}\,t^{n} + a_{1}\,t^{n-1} + \cdots + a_{n-1}\,t + a_{n},
\end{equation}
deren Coefficienten ganze Zahlen sein sollen. Wir wählen ausserdem eine Primzahl $p$ als Modulus, und nehmen an, $a_{0}$ sei durch $p$ nicht theilbar.

Zwei ganze Functionen von $t$, in denen entsprechende Coefficienten nach dem Modul $p$ congruent sind, sollen selbst nach dem Modul $p$ \emph{congruent} genannt werden.

Die Function $f(t)$ heisst nach dem Modul $p$ \emph{reducibel}, wenn es zwei ganze ganzzahlige Functionen $g(t)$, $\psi(t)$ giebt, in deren jeder wenigstens ein von $t$ abhängiges Glied einen durch $p$ nicht theilbaren Coefficienten hat, so dass
\begin{equation}\label{eq:63.2}
\tag{2}
f(t) \equiv g(t)\,\psi(t) \pmod{p}
\end{equation}
ist. In $g(t)$ und $\psi(t)$ kann man alle Glieder, deren Coefficienten durch $p$ theilbar sind, weglassen, und dann muss der Grad von $g(t)\,\psi(t)$ mit dem Grade von $f(t)$ übereinstimmen. Es müssen also die Grade von $g(t)$ und $\psi(t)$ niedriger als $n$ sein. Giebt es keine solchen Functionen $g(t)$, $\psi(t)$, so heisst $f(t)$ nach dem Modul $p$ \emph{irreducibel}.

Eine nach dem Modul $p$ irreducible Function ist gewiss immer im Körper der rationalen Zahlen absolut irreducibel. Das Umgekehrte ist aber nicht nothwendig.

So ist z.~B. die Function $2^{\text{ten}}$ Grades $t^{2} - R$ reducibel nach $p$, wenn $R$ quadratischer Rest von $p$ ist; denn ist $a^{2} \equiv R \pmod{p}$, so ist
\[
t^{2} - R \equiv (t-a)\,(t+a) \pmod{p}.
\]
Dagegen ist $t^{2} - N$, wenn $N$ Nichtrest von $p$ ist, irreducibel, weil sonst $t^{2} - N$ für ein rationales $t$ durch $p$ theilbar werden müsste.

Ein anderes Beispiel bieten die Functionen
\[
t^{2} + t + 1,\quad t^{2} + t + 1,
\]
die für den Modul~2 irreducibel sind. Denn wären sie reducibel, so müsste einer der Factoren linear sein, und es müsste eine ganze Zahl geben, die, für $t$ eingesetzt, diese Functionen zu geraden Zahlen macht. Das aber ist offenbar unmöglich, da beide Functionen für $t = 0$ und $t = 1$ ungerade Zahlen sind.

\newpage
\fancyhead[C]{\small Neunter Abschnitt. \hfill \thepage}

Bedeutet
\begin{equation}\label{eq:63.3}
\tag{3}
P = t^{n} + p_{1}\,t^{n-1} + p_{2}\,t^{n-2} + \cdots + p_{n-1}\,t + p_{n}
\end{equation}
eine nach dem Modul $p$ irreducible Function $n^{\text{ten}}$ Grades, in der wir der Einfachheit halber den Coefficienten der höchsten Potenz $t^{n}$ gleich~1 annehmen, so lässt sich aus jeder anderen ganzen Function $F(t)$ mit ganzzahligen Coefficienten ein Rest $\varrho(t)$ ableiten, dessen Grad niedriger als $n$ ist, indem man (nach \S.~3 des ersten Bandes)
\begin{equation}\label{eq:63.4}
\tag{4}
F(t) = Q\,P + \varrho
\end{equation}
setzt, $\varrho(t)$ hat hierin ganzzahlige Coefficienten, und wir bezeichnen ihre Beziehung zur Function $F(t)$ als eine \emph{Congruenz nach dem Modul}~$P$.

Es heissen also zwei Functionen $F(t)$, $F_{1}(t)$, die denselben Rest $\varrho(t)$ haben, \emph{congruent nach dem Modul}~$P$, was durch die Formel
\[
F(t) \equiv F_{1}(t) \pmod{P}
\]
ausgedrückt wird. Damit ist gleichbedeutend, dass $F(t) - F_{1}(t)$ durch $P$ theilbar ist.

Wenn zwei Functionen $F(t)$ und $F_{1}(t)$ nicht gleichen Rest geben, sondern zwei Reste, die nach dem Modul $p$ congruent sind, so heissen die Functionen $F$, $F_{1}$ \emph{congruent nach dem Doppelmodul} $P$, $p$, und man drückt dies durch eine Formel so aus:
\begin{equation}\label{eq:63.5}
\tag{5}
F(t) \equiv F_{1}(t) \pmod{P, p}.
\end{equation}
Für diese Art der Congruenz gilt der Satz:

\medskip\noindent
\textbf{1.} \emph{Das Product zweier Functionen $F(t)$, $F_{1}(t)$ kann nicht mit Null congruent sein, wenn nicht der eine Factor mit Null congruent ist.}

Der Beweis ergiebt sich aus dem Algorithmus des grössten gemeinschaftlichen Theilers. Ist nämlich $P_{1}$ eine ganze Function von niederigerem Grade als $P$, die nicht nach dem Modul $p$ mit Null congruent ist, so kann man eine Reihe von eben solchen Functionen $P_{2}$, $P_{3}$, $\ldots$ von abnehmendem Grade, und die Quotienten $Q_{1}$, $Q_{2}$, $\ldots$ gleichfalls als ganze Functionen so bestimmen, dass
\begin{equation}\label{eq:63.6}
\tag{6}
\begin{aligned}
P &= Q_{1}\,P_{1} + P_{2},\\
P_{1} &= Q_{2}\,P_{2} + P_{3},\;\ldots\\
P_{\nu-1} &= Q_{\nu-1}\,P_{\nu-1} + P_{\nu} \pmod{p}
\end{aligned}
\end{equation}
wird, und die Reihe dieser Gleichungen bricht ab, wenn $P_{\nu}$ eine Constante (ganze Zahl) geworden ist. Diese Constante $P_{\nu}$ kann aber nicht $\equiv 0 \pmod{p}$ sein; denn sonn liesse sich aus (6) eine Congruenz ableiten:
\[
P \equiv T\,P_{\nu-1} \pmod{p},
\]
in der $T$ eine nicht constante Function von $t$ ist, und $P$ wäre nicht irreducibel nach dem Modul $p$.

Ist nun $Q$ irgend eine ganze Function von $t$, die der Bedingung
\[
Q\,P_{1} \equiv 0 \pmod{P, p}
\]
genügt, so folgt aus (6) durch Multiplication mit $Q$:
\[
Q\,P_{2} \equiv 0,\quad Q\,P_{3} \equiv 0,\;\ldots,\quad Q\,P_{\nu} \equiv 0 \pmod{P, p},
\]
also $Q \equiv 0$. Und wenn wir also für $P_{1}$, $Q$ die Reste der Functionen $F(t)$, $F_{1}(t)$ setzen, so ist hiermit der Satz~1. bewiesen.

Die Reste aller Functionen $F(t)$ nach dem Modul $P$ sind von der Form
\begin{equation}\label{eq:63.7}
\tag{7}
X = x_{0}\,t^{n-1} + x_{1}\,t^{n-2} + \cdots + x_{n-1},
\end{equation}
und wenn man nur die nach dem Modul $p$ incongruenten unter ihnen haben will, so genügt es, die Coefficienten $x_{0}$, $x_{1}$, $\ldots$, $x_{n-1}$ die Reihe der Zahlen $0, 1, 2, \ldots, p-1$ durchlaufen zu lassen.

Es giebt also $p^{n}$ und nicht mehr nach den Moduln $P$, $p$ incongruente Functionen.

% §. 64
\newpage
\sect{64}
\begin{center}
\textbf{Congruenzkörper.}
\end{center}
\label{sec:64}
\fancyhead[C]{\small\S.~64. Congruenzkörper. \hfill 245}

Es ist nur eine andere Ausdrucksweise für die im vorigen Paragraphen durchgeführten Betrachtungen, wenn man mit Galois eine imaginäre Zahl $\varepsilon$ einführt, die der Congruenz
\begin{equation}\label{eq:64.1}
\tag{1}
P(\varepsilon) \equiv 0 \pmod{p}
\end{equation}
genügt. Eine ganze rationale Zahl dieser Art giebt es nicht, sobald $n > 1$ ist, weil sonst $P(t) - P(\varepsilon) = (t - \varepsilon)\,Q$ durch $t - \varepsilon$ theilbar, mithin $P(t) = (t - \varepsilon)\,Q$, und $P(t)$ nicht irreducibel wäre. Gleichwohl kann man, wie man in der gewöhnlichen Zahlenlehre die imaginäre Einheit $i = \sqrt{-1}$ einführt, ein solches Symbol $\varepsilon$ in der Rechnung benutzen. Man rechnet damit, sofern es sich um Addition, Subtraction und Multiplication handelt, nach den Regeln der Buchstabenrechnung, und kann dabei nach Belieben die Gleichung $P(\varepsilon) = 0$ benutzen. Man kann sogar unter $\varepsilon$ geradezu eine Wurzel der Gleichung $P(z) = 0$ im gewöhnlichen Sinne verstehen. Alle Zahlen, auf die man hierbei kommt, lassen sich auf die Form
\begin{equation}\label{eq:64.2}
\tag{2}
\alpha = a_{0}\,\varepsilon^{n-1} + a_{1}\,\varepsilon^{n-2} + \cdots + a_{n-1}
\end{equation}
bringen. Da bei der Rechnung mit solchen Zahlen der Modul $p$ immer festgehalten wird, so wollen wir zwei Zahlen, die nach dem Modul $p$ congruent sind, geradezu als \emph{gleich} bezeichnen. Objecte der Rechnung sind dann eigentlich nicht die einzelnen Zahlen selbst, sondern die aus allen unter einander congruenten Zahlen bestehenden Zahlclassen. Diese Zahlen $\varepsilon$ werden die \emph{Galois'schen Imaginären} genannt. Das zu einem bestimmten $\varepsilon$ gehörige System solcher imaginärer Zahlen bezeichnen wir der Abkürzung wegen mit~$\mathfrak{E}$.

Wir haben dann zunächst den Satz:

\medskip\noindent
\textbf{2.} \emph{Es giebt $p^{n}$ und nicht mehr verschiedene Zahlen in $\mathfrak{E}$.}

Aus dem Satze 1.~aber folgt noch:

\medskip\noindent
\textbf{3.} \emph{Das Product von zwei oder mehr Zahlen aus $\mathfrak{E}$ ist dann und nur dann gleich Null, wenn wenigstens einer der Factoren gleich Null ist.}

Wir erhalten das vollständige Zahlensystem $\mathfrak{E}$, wenn man die rationalen Zahlen $a_{0}$, $a_{1}$, $\ldots$, $a_{n-1}$ in (2) je ein volles Restsystem nach dem Modul $p$ durchlaufen lässt. Jedes System $\mathfrak{E}$ enthält die Reste der natürlichen Zahlen $0, 1, \ldots, p-1$, und für $n = 1$ ist $\mathfrak{E}$ mit diesem Systeme identisch.

Ist $\alpha$ eine feste von Null verschiedene Zahl in $\mathfrak{E}$, so durchläuft $\alpha\,\xi$ zugleich mit $\xi$ das volle System $\mathfrak{E}$. Denn aus 3.~folgt, dass $\alpha\,\xi$ nur dann gleich $\alpha\,\xi'$ sein kann, wenn $\xi = \xi'$ ist. Daraus folgt:

\medskip\noindent
\textbf{4.} \emph{Sind $\alpha$, $\beta$ zwei gegebene Zahlen in $\mathfrak{E}$ und $\alpha$ von Null verschieden, so giebt es eine und nur eine Zahl $\gamma$ in $\mathfrak{E}$, die der Bedingung}
\[
\alpha\,\gamma = \beta
\]
\emph{genügt.}

Damit ist auch die Operation der Division in dem Systeme als erlaubt nachgewiesen. Wir bezeichnen die Zahl $\gamma$, deren Existenz in 4.~ausgesprochen ist, mit
\[
\beta : \alpha\quad\text{oder}\quad \frac{\beta}{\alpha}.
\]

\newpage
\fancyhead[C]{\small\S.~64. Congruenzkörper. \hfill 247}

Man kann das System $\mathfrak{E}$ einen \emph{endlichen Körper} nennen, da es nur eine endliche Anzahl von Zahlen enthält, die das charakteristische Merkmal eines Körpers, nämlich die unbeschränkte Ausführbarkeit der Rechenoperationen, ausgenommen die Division durch Null, aufweisen (Bd.~I, \S.~139)\footnote{Man sehe des Verfassers Abhandlung: Die allgemeinen Grundlagen der Galois'schen Gleichungstheorie. Mathematische Annalen, Bd.~43.}. Wir wollen daher $\mathfrak{E}$ einen \emph{Congruenzkörper} nennen.

Es soll $n$ der \emph{Grad} und $p$ der \emph{Modul} des Körpers $\mathfrak{E}$ heissen.

Ist $\alpha$ eine von Null verschiedene Zahl in $\mathfrak{E}$, so durchläuft das Product $\alpha\,\xi = \eta$ zugleich mit $\xi$ das volle Zahlensystem $\mathfrak{E}$. Schliessen wir $\xi = 0$ aus, so kommt auch $\eta = 0$ nicht vor, und das Product $\prod \xi$ aller $\xi$ stimmt mit dem Producte aller $\eta$ überein und ist von Null verschieden. Durch Multiplication aller Gleichungen $\alpha\,\xi = \eta$ folgt aber
\[
\alpha^{p^{n}-1}\,\prod \xi = \prod \eta,
\]
und nach Abwerfung des gemeinsamen Factors $\prod \xi$ ergiebt sich

\medskip\noindent
\textbf{5.} \emph{der Fermat'sche Satz}:
\begin{equation}\label{eq:64.3}
\tag{3}
\alpha^{p^{n}-1} = 1.
\end{equation}

Multiplicirt man mit $\alpha$, so erhält man diesen Satz in der Form
\begin{equation}\label{eq:64.4}
\tag{4}
\alpha^{p^{n}} = \alpha,
\end{equation}
in der er auch noch für $\alpha = 0$ besteht. Ist $\alpha$ von Null verschieden, so giebt es wegen (3) einen kleinsten positiven Exponenten $e$, für den
\begin{equation}\label{eq:64.5}
\tag{5}
\alpha^{e} = 1
\end{equation}
ist, und für den folglich die Potenzen $1$, $\alpha$, $\alpha^{2}$, $\ldots$, $\alpha^{e-1}$ von einander verschieden sind. Man sagt dann, $\alpha$ \emph{gehört zum Exponenten}~$e$. Dieser Exponent $e$ muss ein Theiler von $p^{n} - 1$ sein. Denn ist $\alpha^{m} = 1$ für irgend einen Exponenten $m$, so setzen wir $m = qe + e'$, worin $q$ eine ganze Zahl und $e' < e$ ist. Dann ist auch $\alpha^{e'} = 1$, und folglich muss $e' = 0$, d.~h. $m$ durch $e$ theilbar sein. Unter diesen Exponenten $m$ findet sich auch $p^{n} - 1$.

Setzen wir also
\[
p^{n} - 1 = ef,
\]
so wird $\alpha^{h}$ immer dann zum Exponenten $e$ gehören, wenn $h$ relativ prim zu $e$ ist. Denn dann und nur dann ist $he$ das kleinste positive, durch $e$ theilbare Vielfache von~$h$.

\newpage
\fancyhead[C]{\small\S.~64. Quadrate und Nichtquadrate. \hfill 249}

Giebt es also überhaupt eine zum Exponenten $e$ gehörige Zahl $\alpha$, so giebt es so viele wie relative Primzahlen zu $e$, die positiv und kleiner als $e$ sind, eine Zahl, die wir schon früher mit $\varphi(e)$ bezeichnet haben, und die, wenn $e$ alle Divisoren von $p^{n} - 1$ durchläuft, der Relation
\begin{equation}\label{eq:64.6}
\tag{6}
\sum \varphi(e) = p^{n} - 1
\end{equation}
genügt (Bd.~I, \S.~132, 133).

Ist $\psi(e)$ die Anzahl der Zahlen $\alpha$, die zum Exponenten $e$ gehören, so ist $\psi(e) = \varphi(e)$ oder $= 0$, und da die Anzahl aller Zahlen $\alpha$ gleich $p^{n} - 1$ ist und jede Zahl $\alpha$ zu einem und nur zu einem Exponenten gehört, so ist auch
\begin{equation}\label{eq:64.7}
\tag{7}
\sum \psi(e) = p^{n} - 1,
\end{equation}
und aus (6) und (7) folgt, dass $\psi(e)$ immer gleich $\varphi(e)$ ist.

Es giebt also $\varphi(p^{n} - 1)$ Zahlen $\gamma$ in $\mathfrak{E}$, die zum Exponenten $p^{n} - 1$ gehören, und die folglich die Eigenschaft haben, dass die Potenzen
\begin{equation}\label{eq:64.8}
\tag{8}
1,\; \gamma,\; \gamma^{2},\; \ldots,\; \gamma^{p^{n}-2}
\end{equation}
alle von einander verschieden sind. Durch diese Reihe ist daher die Gesammtheit der von Null verschiedenen Zahlen in $\mathfrak{E}$ erschöpft.

Solche Zahlen $\gamma$ heissen \emph{primitive Wurzeln} von $\mathfrak{E}$.

Unter den Zahlen in $\mathfrak{E}$ giebt es solche, die als Quadrat einer anderen Zahl in $\mathfrak{E}$ darstellbar sind, die wir \emph{Quadrate} nennen, und andere, bei denen dies nicht zutrifft, die wir \emph{Nichtquadrate} nennen.

Wenn $p = 2$ ist, so ist jede Zahl in $\mathfrak{E}$ ein Quadrat, wie die Formel (4) zeigt.

Wenn aber $p$ ungerade ist, so besteht die eine Hälfte der von Null verschiedenen Zahlen in $\mathfrak{E}$ aus Quadraten, die andere aus Nichtquadraten.

Denn zwei entgegengesetzte Zahlen, wie $+a$ und $-a$, sind bei ungeradem $p$ von einander verschieden, und geben trotzdem dasselbe Quadrat. Wenn man also die sämmtlichen Zahlen von $\mathfrak{E}$ zum Quadrat erhebt, so erhält man höchstens $\frac{1}{2}(p^{n} - 1)$ verschiedene Quadrate. Nun kann andererseits $\beta^{2}$ nur dann gleich $\alpha^{2}$ sein, wenn $\beta = \pm\alpha$ ist; denn aus $\beta^{2} - \alpha^{2} = (\beta - \alpha)\,(\beta + \alpha) = 0$ folgt, dass $\beta - \alpha$ oder $\beta + \alpha$ verschwinden muss. Es giebt also wirklich $\frac{1}{2}(p^{n} - 1)$ Quadrate und ebenso viele Nichtquadrate.

Wenn man die Null mit zu den Quadraten zählt, so erhält man den Satz:

\medskip\noindent
\textbf{6.} \emph{Wenn $p = 2$ ist, so sind alle Zahlen in $\mathfrak{E}$ Quadrate. Ist $p$ ungerade, so giebt es $\frac{1}{2}(p^{n} + 1)$ Quadrate, $\frac{1}{2}(p^{n} - 1)$ Nichtquadrate in $\mathfrak{E}$.}

Bei ungeradem $p$ muss eine primitive Wurzel immer ein Nichtquadrat sein und in der Reihe (8) sind die Zahlen mit geraden Exponenten die Quadrate, die mit ungeraden Exponenten die Nichtquadrate.

Das Product und der Quotient zweier Quadrate ist offenbar wieder ein Quadrat.

Daraus folgt, dass das Product aus einem Nichtquadrat und einem von Null verschiedenen Quadrate ein Nichtquadrat ist. Denn ist das Product $\alpha\beta$ ein Quadrat und der eine Factor $\beta$ ein Quadrat, so muss auch $\alpha = \alpha\beta : \beta$ ein Quadrat sein.

Weiter schliesst man daraus, dass das Product von zwei Nichtquadraten ein Quadrat ist. Denn bedeutet $\beta$ irgend ein Nichtquadrat, so lasse man in dem Producte $\beta\,\xi$ den Factor $\xi$ die sämmtlichen von Null verschiedenen Zahlen von $\mathfrak{E}$ durchlaufen. Das Product durchläuft dann dieselbe Zahlenreihe, und da für alle quadratischen $\xi$ das Product $\beta\,\xi$ Nichtquadrat ist, so muss es für die nichtquadratischen $\xi$ Quadrat sein.

Alles dies ergiebt sich auch sehr einfach aus der Darstellung (8) der Zahlen von $\mathfrak{E}$ durch eine primitive Wurzel.

Wir schliessen diese allgemeinen Betrachtungen mit dem Satze, den wir später brauchen werden:

\medskip\noindent
\textbf{7.} \emph{Jede nicht quadratische Zahl ist die Summe von zwei Quadraten.}

Hierbei ist $p$ als ungerade vorauszusetzen, da es nur dann nichtquadratische Zahlen giebt.

Es lässt sich dann zunächst jede Zahl $\beta$ als Differenz zweier Quadrate darstellen, wie die Identität
\[
(\beta + \tfrac{1}{4})^{2} - (\beta - \tfrac{1}{4})^{2} = \beta
\]
zeigt. Ist nun $-1$ ein Quadrat in $\mathfrak{E}$, so setze man
\[
(\beta + \tfrac{1}{4})^{2} = \sigma,\quad -(\beta - \tfrac{1}{4})^{2} = \tau,\qquad
\beta = \sigma + \tau.
\]
Ist aber $-1$ und also auch $\beta - 1$ Nichtquadrat, so suche man in der Reihe der natürlichen Zahlen $1, 2, \ldots, p-1$, deren erstes Glied ein Quadrat und deren letztes ein Nichtquadrat ist, ein quadratisches Glied auf, auf welches ein Nichtquadrat folgt. Ist dann also $\alpha = \varrho^{2}$ ein Quadrat und $\varrho^{2} + 1$ ein Nichtquadrat, so ist, wenn $\beta$ ein Nichtquadrat ist, $\beta : (\varrho^{2} + 1) = \tau^{2}$ ein Quadrat, und daraus folgt:
\[
\beta = \tau^{2}\,\varrho^{2} + \tau^{2} = \sigma^{2} + \tau^{2}.\qquad\text{w.~z.~b.~w.}
\]



% §. 65
\newpage
\sect{65}
\begin{center}
\textbf{Congruenzgruppen im Körper $\mathfrak{E}$.}
\end{center}
\label{sec:65}
\fancyhead[C]{\small\S.~65. Congruenzgruppen. \hfill 251}

In jedem Congruenzkörper $\mathfrak{E}$ kann man eine endliche Gruppe von linearen Substitutionen bilden, wenn man in
\begin{equation}\label{eq:65.1}
\tag{1}
A = \begin{pmatrix}\alpha&\beta\\\gamma&\delta\end{pmatrix}
\end{equation}
die Elemente $\alpha$, $\beta$, $\gamma$, $\delta$ alle Zahlen von $\mathfrak{E}$ durchlaufen lässt, und wenn man je zwei dieser Substitutionen nach der Vorschrift
\begin{equation}\label{eq:65.2}
\tag{2}
\begin{pmatrix}\alpha&\beta\\\gamma&\delta\end{pmatrix}\,
\begin{pmatrix}\alpha'&\beta'\\\gamma'&\delta'\end{pmatrix}
=
\begin{pmatrix}
\alpha\alpha' + \beta\gamma',& \alpha\beta' + \beta\delta'\\
\gamma\alpha' + \delta\gamma',& \gamma\beta' + \delta\delta'
\end{pmatrix}
\end{equation}
componirt. Der Grad dieser Gruppe ist $p^{4n}$, wenn $p$ der Modul und $n$ der Grad von $\mathfrak{E}$ ist.

Darin ist aber eine Gruppe niedrigeren Grades enthalten. Nennen wir $\alpha\delta - \beta\gamma$ die \emph{Determinante} der Substitution (1), so folgt aus (2), dass die Determinante der aus $A$ und $A'$ componirten Substitution $AA'$ gleich dem Producte der Determinanten von $A$ und $A'$ ist. Wenn wir daher den Zahlen $\alpha$, $\beta$, $\gamma$, $\delta$ die Bedingung auferlegen, dass
\begin{equation}\label{eq:65.3}
\tag{3}
\alpha\delta - \beta\gamma = 1
\end{equation}
sein soll, so bilden auch diese Elemente $A$ eine Gruppe. Der Grad dieser Gruppe ist gleich der Anzahl der Lösungen von~(3).

Um diese Anzahl zu finden, bemerken wir zunächst, dass wir für $\alpha$, $\beta$ irgend zwei Zahlen aus $\mathfrak{E}$ setzen können, die nicht beide gleich Null sind. Denn ist etwa $\alpha$ von Null verschieden, so kann man $\gamma = 0$, $\delta = 1 : \alpha$ setzen und hat so die Bedingung (3) erfüllt. Die Anzahl der brauchbaren Zahlenpaare $\alpha$, $\beta$ ist also $p^{2n} - 1$. Hat man aber zu einem Zahlenpaare $\alpha$, $\beta$ eine Lösung $\gamma_{0}$, $\delta_{0}$ von (3) gefunden, so müssen alle übrigen der Bedingung
\[
\alpha\,(\delta - \delta_{0}) = \beta\,(\gamma - \gamma_{0})
\]
genügen, und wenn man also $\gamma - \gamma_{0} = \alpha\,\xi$ setzt (falls $\alpha$ von Null verschieden ist), so folgt $\delta - \delta_{0} = \beta\,\xi$, also
\[
\gamma = \gamma_{0} + \alpha\,\xi,\quad \delta = \delta_{0} + \beta\,\xi,
\]
umgekehrt genügt auch jedes in dieser Form enthaltene Zahlenpaar $\gamma$, $\delta$. Dies gilt offenbar auch noch in dem ausgenommenen Falle $\alpha = 0$. Da wir nun $p^{n}$ verschiedene Zahlen für $\xi$ setzen können, so ergiebt sich der Grad unserer Gruppe gleich
\begin{equation}\label{eq:65.4}
\tag{4}
p^{n}\,(p^{2n} - 1).
\end{equation}

Ist $p$ ungerade, so können wir eine noch kleinere Gruppe ableiten:

Ist nämlich $p = 2$, so ist eine Zahl $\alpha$ von $-\alpha$ nicht verschieden, es sind also auch die Elemente
\begin{equation}\label{eq:65.5}
\tag{5}
\begin{pmatrix}\alpha&\beta\\\gamma&\delta\end{pmatrix}
\quad\text{und}\quad
\begin{pmatrix}-\alpha&-\beta\\-\gamma&-\delta\end{pmatrix}
\end{equation}
mit einander identisch. Wenn aber $p$ ungerade ist, so sind die beiden Elemente (5) von einander verschieden, und das ganze System dieser Substitutionen zerfällt in Paare von der Form (5).

Wenn wir zwei solche Paare nehmen und componiren je ein Element des einen Paares mit je einem Element des anderen nach der Regel (2), so erhalten wir nur zwei verschiedene Elemente, die wieder ein solches Paar bilden. Diese Paare können also selbst als Elemente einer neuen Gruppe vom Grade
\begin{equation}\label{eq:65.6}
\tag{6}
\frac{p^{n}\,(p^{2n} - 1)}{2}
\end{equation}
aufgefasst werden. Wir können uns auch so ausdrücken, dass zwei Elemente $A$, deren entsprechende Zahlen nur durch das Vorzeichen unterschieden sind, als nicht verschieden anzusehen sind.

Die so definirte Gruppe, deren Grad für $p = 2$ durch (4) und für ein ungerades $p$ durch (6) ausgedrückt ist, wollen wir die zum Körper $\mathfrak{E}$ \emph{gehörige Congruenzgruppe} nennen und mit $\mathfrak{E}$ bezeichnen.

Für die genauere Untersuchung dieser Gruppe ist es von Wichtigkeit, ein System erzeugender Substitutionen zu kennen. Stellen wir nach \S.~64,~(2) die Zahlen von $\mathfrak{E}$ in der Form dar:
\begin{equation}\label{eq:65.7}
\tag{7}
\xi = z_{0} + z_{1}\,\varepsilon + z_{2}\,\varepsilon^{2} + \cdots + z_{n-1}\,\varepsilon^{n-1},
\end{equation}
worin die $z_{i}$ rationale, nach dem Modul $p$ genommene Zahlen bedeuten, so erhalten wir ein System erzeugender Substitutionen in der Form:
\begin{equation}\label{eq:65.8}
\tag{8}
A_{1} = \begin{pmatrix}1&1\\0&1\end{pmatrix},\quad
B_{\xi} = \begin{pmatrix}1&0\\\xi&1\end{pmatrix},\quad
A_{\eta} = \begin{pmatrix}\eta&0\\0&\eta^{-1}\end{pmatrix},\quad
C = \begin{pmatrix}0&-1\\1&0\end{pmatrix} = I.
\end{equation}

Um dies nachzuweisen, bemerken wir, dass
\[
\begin{pmatrix}1&\xi\\0&1\end{pmatrix}\,
\begin{pmatrix}1&\xi'\\0&1\end{pmatrix}
= \begin{pmatrix}1&\xi + \xi'\\0&1\end{pmatrix}
\]
ist; und durch wiederholte Anwendung hiervon, folgt, wenn
\[
\xi = z_{0} + z_{1}\,\varepsilon + \cdots + z_{n-1}\,\varepsilon^{n-1}
\]
mit ganzen rationalen Coefficienten $z$ eine beliebige Zahl in $\mathfrak{E}$ ist,
\begin{equation}\label{eq:65.9}
\tag{9}
\begin{pmatrix}1&\xi\\0&1\end{pmatrix}
= B_{z_{0}}\,B_{z_{1}\varepsilon}\,\cdots\,B_{z_{n-1}\varepsilon^{n-1}}.
\end{equation}
Ferner folgt durch Zusammensetzung mit $A_{\eta}$:
\begin{equation}\label{eq:65.10}
\tag{10}
\begin{pmatrix}\eta&\eta\,\xi\\0&\eta^{-1}\end{pmatrix}
= A_{\eta}\,\begin{pmatrix}1&\xi\\0&1\end{pmatrix},\qquad
\begin{pmatrix}\eta&0\\\xi&\eta^{-1}\end{pmatrix}
= \begin{pmatrix}1&0\\\xi\,\eta^{-1}&1\end{pmatrix}\,A_{\eta}.
\end{equation}

Daraus geht hervor, dass man aus den Elementen (8) alle Substitutionen von der Form
\[
\begin{pmatrix}\eta&\xi\\0&\eta^{-1}\end{pmatrix},\quad
\begin{pmatrix}\eta&0\\\xi&\eta^{-1}\end{pmatrix}
\]
zusammensetzen kann, worin $\xi$, $\eta$ beliebige Zahlen in $\mathfrak{E}$ sind. Ferner ist
\begin{equation}\label{eq:65.11}
\tag{11}
\begin{pmatrix}\eta&\xi\\0&\eta^{-1}\end{pmatrix}\,
\begin{pmatrix}0&-1\\1&0\end{pmatrix}
= \begin{pmatrix}\xi&-\eta\\\eta^{-1}&0\end{pmatrix}.
\end{equation}
Nimmt man $\xi$ von Null verschieden an, so kann man $\alpha$ und $\gamma$ so wählen, dass $\alpha\xi + 1$, $-\beta\gamma - 1$ beliebige Zahlen werden, und nach (11) kann man also alle Substitutionen $\begin{pmatrix}\alpha&\beta\\\gamma&\delta\end{pmatrix}$, worin $\delta$ von Null verschieden ist, aus $A_{1}$, $B_{\xi}$, $A_{\eta}$, $I$ zusammensetzen. Die Beschränkung, dass $\delta$ von Null verschieden sei, kann aber nach der Formel
\[
\begin{pmatrix}\alpha&\beta\\0&0\end{pmatrix}
= \begin{pmatrix}0&-1\\1&0\end{pmatrix}\,
\begin{pmatrix}0&-1\\1&0\end{pmatrix}\,
\begin{pmatrix}0&1\\-1&\alpha\end{pmatrix}
\]
fallen gelassen werden.

\newpage
\fancyhead[C]{\small\S.~65. Congruenzgruppen. \hfill 253}

Die Gruppe $\mathfrak{E}$ enthält einen Theiler vom Index $p^{n} + 1$, der aus allen Substitutionen der Form
\[
\begin{pmatrix}\xi&\eta\\0&\xi^{-1}\end{pmatrix}
\]
besteht, worin $\xi$ jede beliebige, $\eta$ jede von Null verschiedene Zahl in $\mathfrak{E}$ bedeutet, und $\delta = \alpha^{-1}$ ist.

Diesem Theiler entspricht eine der Gruppe $\mathfrak{E}$ isomorphe Permutationsgruppe von $p^{n} + 1$ Ziffern, die man so bilden kann:

Der lineare Ausdruck
\begin{equation}\label{eq:65.12}
\tag{12}
\eta = \frac{\alpha\,\xi + \beta}{\gamma\,\xi + \delta}
\end{equation}
giebt für jede Zahl $\xi$ aus $\mathfrak{E}$ eine entsprechende Zahl $\eta$, ausgenommen, wenn $\gamma\,\xi + \delta = 0$ ist. Ebenso giebt es zu jedem $\eta$ ein bestimmtes $\xi$, ausgenommen für $\gamma\,\eta - \alpha = 0$. Um diese Ausnahme zu beseitigen, genügt es, das System $\mathfrak{E}$ durch Hinzufügung eines Elementes, das wir \emph{„Unendlich"} nennen und mit $\infty$ bezeichnen, zu erweitern. Mit diesem Zeichen wird nach den folgenden Regeln gerechnet, worin $a$ irgend ein Element des erweiterten Systemes $\mathfrak{E}$ bedeutet:
\[
\begin{aligned}
a + \infty &= \infty, & a - \infty &= \infty,\quad &&\text{ausser für } a = \infty\\
a \cdot \infty &= \infty, & &&&&\text{für } a \neq 0\\
a : 0 &= \infty, &&&&&&\text{für } a \neq 0\\
a : \infty &= 0, &&&&&&\text{für } a \neq \infty.
\end{aligned}
\]
Dann führt das Ergebniss jeder Rechnung mit den vier Species immer zu einer bestimmten Zahl des Systems $\mathfrak{E}$ und nur die Verbindungen $\infty + \infty$, $0 \cdot \infty$, $\infty : \infty$, $0 : 0$ bleiben ohne Bedeutung.

Nach (12) entspricht dann der Zahl $\xi = -\delta : \gamma$ die Zahl $\eta = \infty$, und der Zahl $\xi = \infty$ die Zahl $\eta = \alpha : \gamma$, und jeder Substitution $A$ in $\mathfrak{E}$ entspricht eine bestimmte Permutation der $p^{n} + 1$ Elemente des erweiterten Systemes $\mathfrak{E}$. Man sieht leicht, dass nur die identische Substitution alle diese Elemente ungeändert lässt, und dass folglich auch zwei verschiedene Substitutionen aus $\mathfrak{E}$ immer verschiedene Permutationen hervorrufen. Der Isomorphismus ist also einstufig.

% §. 66
\newpage
\sect{66}
\begin{center}
\textbf{Einfachheit der Gruppe $\mathfrak{E}$.}
\end{center}
\label{sec:66}
\fancyhead[C]{\small\S.~66. Einfachheit der Gruppe $\mathfrak{E}$. \hfill 255}

Die erste Frage bei einer eingehenderen Untersuchung der Gruppe $\mathfrak{E}$ ist die nach einem etwa vorhandenen Normaltheiler. Es lässt sich beweisen, dass ein solcher Normaltheiler, von zwei ganz einfachen Ausnahmen abgesehen, nicht vorhanden ist.

Nehmen wir an, es sei $G$ ein Normaltheiler von $\mathfrak{E}$, der wenigstens eine von der Identität verschiedene Substitution
\begin{equation}\label{eq:66.1}
\tag{1}
A = \begin{pmatrix}\alpha&\beta\\\gamma&\delta\end{pmatrix}
\end{equation}
enthält. Nach dem Begriffe des Normaltheilers ist dann, wenn $T$ ein beliebiges Element in $\mathfrak{E}$ ist,
\begin{equation}\label{eq:66.2}
\tag{2}
T\,A\,T^{-1}
\end{equation}
auch in $G$ enthalten, und es ist zu zeigen, dass man auf diese Weise und durch Zusammensetzung solcher Substitutionen alle Elemente von $\mathfrak{E}$ ableiten kann, woraus dann folgt, dass $G$ mit $\mathfrak{E}$ identisch sein muss. Es genügt aber dazu, das Vorkommen der erzeugenden Substitutionen $A_{1}$, $B_{\xi}$, $I$ [\S.~65,~(8)] in $G$ nachzuweisen.

Wir gehen dazu schrittweise vor.

\medskip\noindent
\textbf{1)} In $G$ kommt eine Substitution von der Form
\[
\begin{pmatrix}0&\beta\\\gamma&\delta\end{pmatrix}
\]
vor. Um dies zu zeigen, bilden wir nach (2):
\begin{equation}\label{eq:66.3}
\tag{3}
\begin{pmatrix}1&\xi\\0&1\end{pmatrix}\,
\begin{pmatrix}\alpha&\beta\\\gamma&\delta\end{pmatrix}\,
\begin{pmatrix}1&-\xi\\0&1\end{pmatrix}
=
\begin{pmatrix}
\alpha + \gamma\xi,& -\alpha\xi + \beta - \gamma\xi^{2} + \delta\xi\\
\gamma,& -\gamma\xi + \delta
\end{pmatrix}.
\end{equation}
Wenn nun $\gamma$ von Null verschieden ist, so kann man aus $\beta + \delta\,\xi = 0$ $\xi$ bestimmen und hat das Ziel erreicht. Ist aber $\gamma = 0$, so ist nicht gleichzeitig $\delta - \alpha = 0$, $\beta = 0$, weil sonst $A$ die identische Substitution wäre; man kann daher $\xi$ so annehmen, dass $(\delta - \alpha)\,\xi - \beta$ von Null verschieden ist. Dann hat man durch (3) eine Substitution $A$ in $G$ gebildet, in der $\beta$ nicht Null ist, und auf die man die Formel (3) nun nochmals anwenden kann, um $\alpha$ zu Null zu machen.

\medskip\noindent
\textbf{2)} In $G$ kommt eine Substitution von der Form
\[
\begin{pmatrix}1&\beta\\0&1\end{pmatrix}
\]
vor. Um dies nachzuweisen, bilden wir nach (2):
\begin{equation}\label{eq:66.4}
\tag{4}
\begin{pmatrix}\lambda&\mu\\\nu&\varrho\end{pmatrix}\,
\begin{pmatrix}0&\beta\\\gamma&\delta\end{pmatrix}\,
\begin{pmatrix}\varrho&-\mu\\-\nu&\lambda\end{pmatrix}
=
\begin{pmatrix}
-\lambda\nu\delta + \lambda\varrho\beta - \mu\nu\gamma,&
\lambda^{2}\beta - \lambda\mu\delta - \lambda\mu\gamma + \mu^{2}\gamma\\
-\nu^{2}\beta + \nu\varrho\gamma + \nu\varrho\delta - \varrho^{2}\gamma,&
\nu\mu\beta - \nu\mu\delta + \varrho\lambda\gamma - \varrho^{2}\delta
\end{pmatrix},
\end{equation}
und es ist also nur zu zeigen, dass man $\lambda$, $\mu$, $\nu$, $\varrho$ so bestimmen kann, dass die drei Gleichungen:
\[
\begin{aligned}
\lambda\varrho - \mu\nu &= 1\\
-\lambda\nu\delta + \lambda\varrho\beta - \mu\nu\gamma &= 0\\
-\nu^{2}\beta + \nu\varrho\gamma + \nu\varrho\delta - \varrho^{2}\gamma &= 1
\end{aligned}
\]
befriedigt sind. Aus den beiden letzten Gleichungen folgt aber, wenn man sie zuerst mit $\lambda$, $\mu$, dann mit $\nu$, $\varrho$ multiplicirt und jedesmal addirt, mit Benutzung der ersten Gleichung:
\begin{equation}\label{eq:66.5}
\tag{5}
\lambda\beta + \mu\delta = \varrho,\quad \nu\beta = \mu,
\end{equation}
und wenn man diese Werthe von $\varrho$, $\nu$ in die erste Gleichung einsetzt:
\begin{equation}\label{eq:66.6}
\tag{6}
\lambda^{2}\beta + \lambda\mu\delta - \mu^{2}\gamma = 1.
\end{equation}
Diese Gleichung kann aber immer befriedigt werden. Denn wegen $\beta\gamma \not\equiv 0$ kann keine der beiden Zahlen $\beta$, $\gamma$ verschwinden. Ist nun $\beta$ ein Quadrat, so können wir $\mu = 0$ setzen und $\lambda$ aus $\lambda^{2}\beta = 1$ bestimmen. Dies ist immer möglich bei $p = 2$. Ist aber $p$ ungerade, so erhalten wir, wenn $\beta$ ein Nichtquadrat ist, aus (6):
\begin{equation}\label{eq:66.7}
\tag{7}
(2\lambda\beta + \mu\delta)^{2} + \mu\,(4\beta - \delta^{2}) = 4\beta.
\end{equation}
Ist nun zunächst $4\beta - \delta^{2}$ Nichtquadrat, so kann man $\mu$ aus $\mu\,(4\beta - \delta^{2}) = 4$ und dann $\lambda$ aus $2\lambda\beta + \mu\delta = 0$ bestimmen. Ist aber $4\beta - \delta^{2}$ Quadrat, so zerlege man $4\beta$ nach \S.~64,~7. in die Summe aus zwei Quadraten $4\beta = \sigma^{2} + \tau^{2}$ und bestimme $\mu$ und $\lambda$ aus
\[
\mu\,(4\beta - \delta^{2}) = \sigma^{2},\quad 2\lambda\beta + \mu\delta = \tau.
\]
Hat man so $\lambda$ und $\mu$ ermittelt, so erhält man $\varrho$ und $\nu$ aus (5), und dadurch sind die Gleichungen (6) thatsächlich befriedigt. Damit ist das Ziel erreicht.

\medskip\noindent
\textbf{3)} Die Gruppe $G$ enthält die Substitution
\[
\begin{pmatrix}1&1\\0&1\end{pmatrix}
\]
mit alleiniger Ausnahme des Falles $n = 1$, $p = 2$. Zunächst haben wir in $G$ die Substitutionen
\[
\begin{pmatrix}1&\beta\\0&1\end{pmatrix},\quad
\begin{pmatrix}1&\beta'\\0&1\end{pmatrix},\quad
\begin{pmatrix}1&\beta\\0&1\end{pmatrix}\,
\begin{pmatrix}1&\beta'\\0&1\end{pmatrix}
= \begin{pmatrix}1&\beta + \beta'\\0&1\end{pmatrix},
\]
und mithin
\[
\begin{pmatrix}1&m\,\beta\\0&1\end{pmatrix}
= \begin{pmatrix}1&\beta\\0&1\end{pmatrix}^{m}
\]
also auch
\[
\begin{pmatrix}1&1\\0&1\end{pmatrix}
\]
für jede ungerade ganze Zahl $m$; wenn also $p$ ungerade ist, so können wir $m = p$ annehmen und erhalten $A_{1}$. Ist aber $p = 2$, so haben wir, wenn wir mit $\varrho$ eine noch zu bestimmende Zahl in $\mathfrak{E}$ bezeichnen, folgende Substitutionen in $G$:
\[
\begin{pmatrix}1&\varrho^{2}\\0&1\end{pmatrix},\quad
\begin{pmatrix}1&\varrho\\0&1\end{pmatrix},\quad
\begin{pmatrix}1&\varrho^{2}\\0&1\end{pmatrix}\,
\begin{pmatrix}1&\varrho\\0&1\end{pmatrix}
= \begin{pmatrix}1&\varrho^{2} + \varrho\\0&1\end{pmatrix}.
\]
Da hier nun jede Zahl ein Quadrat ist, so kann man $\varrho^{2}$ auch durch $\varrho$ ersetzen, und findet also, dass in $G$ die Substitution
\[
\begin{pmatrix}1&\varrho\\0&1\end{pmatrix}
\]
für jedes beliebige von Null verschiedene $\varrho$, und folglich auch die entgegengesetzte Substitution
\[
\begin{pmatrix}1&-\varrho\\0&1\end{pmatrix}
\]
vorkommen muss. Bedeuten also $\varrho$, $\sigma$ zwei von Null verschiedene Zahlen, so haben wir in $G$ auch
\[
\begin{pmatrix}1&\varrho\\0&1\end{pmatrix}\,
\begin{pmatrix}1&\sigma\\0&1\end{pmatrix}\,
\begin{pmatrix}1&-\varrho\\0&1\end{pmatrix}
= \begin{pmatrix}1&\sigma\\0&1\end{pmatrix}.
\]

\newpage
\fancyhead[C]{\small\S.~66. Einfachheit der Gruppe $\mathfrak{E}$. \hfill 257}

Wenn nun $n$ grösser als 1 ist, so kann man $\varrho$ und $1 + \varrho$ von Null verschieden annehmen und dann $\sigma = -1 : (1 + \varrho)$ setzen, wodurch die Substitution (9) in $A_{1}$ übergeht.

Ist aber $n = 1$ und $p = 2$, so tritt der erste Ausnahmefall ein; denn dann ist immer eine der beiden Zahlen $\varrho$ und $1 + \varrho$ gleich Null. Dieser Fall führt auf eine Gruppe $6^{\text{ten}}$ Grades
\[
\mathfrak{E} =
\begin{pmatrix}1&0\\0&1\end{pmatrix},\;
\begin{pmatrix}1&1\\0&1\end{pmatrix},\;
\begin{pmatrix}1&0\\1&1\end{pmatrix},\;
\begin{pmatrix}0&1\\1&0\end{pmatrix},\;
\begin{pmatrix}1&1\\1&0\end{pmatrix},\;
\begin{pmatrix}0&1\\1&1\end{pmatrix},
\]
in der ein Normaltheiler $3^{\text{ten}}$ Grades enthalten ist:
\[
\mathfrak{G} =
\begin{pmatrix}1&0\\0&1\end{pmatrix},\;
\begin{pmatrix}1&1\\0&1\end{pmatrix},\;
\begin{pmatrix}1&0\\1&1\end{pmatrix}.
\]
Von diesem Ausnahmefalle sehen wir jetzt ab. Wir haben dann weiter:

\medskip\noindent
\textbf{4)} Die Gruppe $G$ enthält jede Substitution von der Form
\[
\begin{pmatrix}1&\xi\\0&1\end{pmatrix},
\]
worin $\xi$ eine beliebige Zahl in $\mathfrak{E}$ ist, ausgenommen in dem Falle $n = 1$, $p = 3$.

Denn nach 3) enthält $G$ die Substitution
\[
\begin{pmatrix}\varrho&0\\0&\varrho^{-1}\end{pmatrix}\,
\begin{pmatrix}1&\beta\\0&1\end{pmatrix}\,
\begin{pmatrix}\varrho^{-1}&0\\0&\varrho\end{pmatrix}
= \begin{pmatrix}1&\varrho^{2}\,\beta\\0&1\end{pmatrix},
\]
wenn $\varrho$ eine beliebige von Null verschiedene Zahl in $\mathfrak{E}$ ist. Daraus folgt, dass auch
\[
\begin{pmatrix}\varrho&0\\0&\varrho^{-1}\end{pmatrix}\,
\begin{pmatrix}1&\beta\\0&1\end{pmatrix}\,
\begin{pmatrix}\varrho&0\\0&\varrho^{-1}\end{pmatrix}\,
\begin{pmatrix}1&-\beta\\0&1\end{pmatrix}
= \begin{pmatrix}1&(\varrho^{2} - 1)\,\beta\\0&1\end{pmatrix}
\]
für jedes beliebige $\beta$ in $G$ vorkommt. Kann man nun $\varrho$ von Null verschieden so annehmen, dass $\varrho^{2} - 1$ nicht verschwindet, so kann man für jedes $\xi$ die Zahl $\beta$ aus $\beta\,(\varrho^{2} - 1) = \xi$ bestimmen, und erhält also aus (11) jede Substitution der Form (10). Um die Möglichkeit hiervon zu beurtheilen, nehmen wir eine primitive Wurzel $\gamma$ von $\mathfrak{E}$ an und setzen $\varrho = \gamma^{h}$.

Kann man den Exponenten $h$ so annehmen, dass $2h$ nicht durch $p^{n} - 1$ theilbar ist, so genügt $\varrho$ unserer Forderung. Dies ist immer möglich, wenn $p = 2$, $n > 1$ ist, und wenn $p^{n} - 1 > 4$ ist, und es bleiben also nur die beiden Fälle $n = 1$, $p = 3$ und $n = 1$, $p = 5$ noch zweifelhaft.

\newpage
\fancyhead[C]{\small\S.~66. Einfachheit der Gruppe $\mathfrak{E}$. \hfill 259}

In dem Falle $n = 1$, $p = 5$ haben wir aber in $G$ nach (2) und 3) die Substitution
\[
\begin{pmatrix}3&2\\2&1\end{pmatrix}\,
\begin{pmatrix}1&1\\0&1\end{pmatrix}\,
\begin{pmatrix}-1&2\\2&-3\end{pmatrix}
= \begin{pmatrix}1&3\\0&1\end{pmatrix},
\]
und folglich auch
\[
\begin{pmatrix}1&3\\0&1\end{pmatrix}\,
\begin{pmatrix}1&1\\0&1\end{pmatrix}
= \begin{pmatrix}1&-1\\0&1\end{pmatrix},
\]
also, da $3$ und folglich auch $-2\,\xi$ beliebig ist, wieder die Substitutionen (10) und damit also alle Erzeugenden der Gruppe $\mathfrak{E}$, und also ist in allen diesen Fällen $G$ mit $\mathfrak{E}$ identisch und folglich die Gruppe $\mathfrak{E}$ \emph{einfach}.

Der Fall $n = 1$, $p = 3$ bietet aber die zweite wirkliche Ausnahme. In diesem Falle ist die Gruppe $\mathfrak{E}$ vom $12^{\text{ten}}$ Grade und sie hat den Normaltheiler $4^{\text{ten}}$ Grades
\[
\mathfrak{G} =
\begin{pmatrix}1&0\\0&1\end{pmatrix},\;
\begin{pmatrix}1&1\\0&1\end{pmatrix},\;
\begin{pmatrix}1&-1\\0&1\end{pmatrix},\;
\begin{pmatrix}-1&0\\0&-1\end{pmatrix}.
\]
Um sich davon zu überzeugen, braucht man nur die beiden Nebengruppen:
\[
\begin{pmatrix}0&-1\\1&0\end{pmatrix}\,\mathfrak{G}
= \begin{pmatrix}0&-1\\1&0\end{pmatrix},\;
\begin{pmatrix}1&-1\\1&0\end{pmatrix},\;
\begin{pmatrix}1&0\\-1&1\end{pmatrix},\;
\begin{pmatrix}0&1\\-1&0\end{pmatrix}
\]
\[
\begin{pmatrix}1&1\\0&1\end{pmatrix}\,\mathfrak{G}
= \begin{pmatrix}1&1\\0&1\end{pmatrix},\;
\begin{pmatrix}1&-1\\0&1\end{pmatrix},\;
\begin{pmatrix}-1&0\\0&-1\end{pmatrix},\;
\begin{pmatrix}0&1\\1&0\end{pmatrix}
\]
mit $\mathfrak{G}$, $\begin{pmatrix}1&0\\0&1\end{pmatrix}$, $\begin{pmatrix}1&1\\0&1\end{pmatrix}$ zu vergleichen.

Damit ist also allgemein bewiesen, dass, von den beiden Fällen $n = 1$, $p = 2$ und $n = 1$, $p = 3$ abgesehen, die Gruppe $\mathfrak{E}$ einfach ist. Für die ersten Fälle erhält man für die Grade $g$ dieser einfachen Gruppen:
\[
\begin{array}{ll|ll}
n = 1,\; p = 5, & g = 60 & n = 2,\; p = 2, & g = 60\\
p = 7, & g = 168 & p = 3, & g = 360\\
p = 11, & g = 660 & p = 5, & g = 7800\\
p = 13, & g = 1092 & n = 3,\; p = 2, & g = 504\\[4pt]
& & p = 3, & g = 9828\\
& & n = 4,\; p = 2, & g = 4080
\end{array}
\]
\begin{flushright}
\footnotesize
1) Siehe Bulletin of the New York math.\ Society. October 1893. In\\
dem Berichte über den mathematischen Congress in Chicago finden sich\\
zwei Notizen über diese Gruppen von Cole und Moore.
\end{flushright}



% §. 67
\newpage
\sect{67}
\begin{center}
\textbf{Congruenzkörper zweiten Grades.}
\end{center}
\label{sec:67}
\fancyhead[C]{\small\S.~67. Congruenzkörper zweiten Grades. \hfill 259}

In jedem Congruenzkörper $\mathfrak{E}$, für den Modul $p$ vom $n^{\text{ten}}$ Grade ist der Congruenzkörper ersten Grades $\mathfrak{E}_{1,p}$, der aus den nach dem Modul $p$ genommenen rationalen Zahlen besteht, als Theiler enthalten. Daher ist auch in der Gruppe der linearen Substitutionen $\mathfrak{E}_{n,p}$ eine Gruppe als Theiler enthalten, die aus allen Substitutionen der Form
\begin{equation}\label{eq:67.1}
\tag{1}
A = \begin{pmatrix}a&b\\c&d\end{pmatrix},\qquad ad - bc = 1
\end{equation}
besteht, worin $a$, $b$, $c$, $d$ nach dem Modul $p$ genommene ganze rationale Zahlen sind.

Diese Gruppe, die als die Gruppe der Modulargleichungen in der Theorie der elliptischen Functionen auftritt\footnote{Man vergleiche des Verfassers Buch: \,Elliptische Functionen und algebraische Zahlen\,. Braunschweig 1891.}, ist das eigentliche Ziel unserer Betrachtungen. Es bietet aber für die Untersuchung dieser Gruppe, und besonders für die Aufsuchung ihrer Theiler, den grössten Vortheil, diese Gruppe nicht selbstständig für sich zu betrachten, sondern als Theiler einer Gruppe $\mathfrak{E}_{2,p}$. Bezeichnen wir nämlich mit $N$ irgend einen feststehenden quadratischen Nichtrest von $p$, so ist, wie schon oben bemerkt, $t^{2} - N$ nach dem Modul $p$ irreducibel, und wir erhalten einen Congruenzkörper $2^{\text{ten}}$ Grades, wenn wir eine Galois'sche imaginäre Zahl durch die Gleichung
\begin{equation}\label{eq:67.2}
\tag{2}
\varepsilon^{2} = N
\end{equation}
einführen.

Es soll der Deutlichkeit wegen für die nächsten Betrachtungen festgesetzt sein, dass die kleinen lateinischen Buchstaben ganze rationale Zahlen (nach dem Modul $p$ genommen), die wir hier auch \emph{reelle} Zahlen nennen, die griechischen Buchstaben Zahlen des Körpers $\mathfrak{E}_{2}$ bedeuten sollen.

Dieser Körper $\mathfrak{E}_{2}$ hat die für uns sehr wichtige Eigenschaft, dass in ihm alle reellen Zahlen Quadrate sind.

Wenn nämlich $a$ quadratischer Rest von $p$ ist, so können wir $a = \alpha^{2}$ setzen, und wenn $b$ quadratischer Nichtrest ist, so kann $\chi^{2} = bN$ befriedigt werden, also $b = (\chi\,\varepsilon^{-1})^{2}$.

\newpage
\fancyhead[C]{\small\S.~67. Congruenzkörper zweiten Grades. \hfill 261}

Eine Zahl von der Form $b\,\varepsilon$ soll eine \emph{rein imaginäre} Zahl heissen, zwei Zahlen der Form $a + b\,\varepsilon$, $a - b\,\varepsilon$ \emph{conjugirt imaginäre} Zahlen. Die Uebertragung dieser Bezeichnungen aus der Theorie der gewöhnlichen complexen Zahlen auf die Zahlen des Körpers $\mathfrak{E}_{2,p}$ rechtfertigt sich durch die Uebereinstimmung der Rechenregeln und kann zu keinem Missverständniss führen, da in diesen Betrachtungen von den gewöhnlichen complexen Zahlen nicht die Rede ist.

Weil immer, wenn $N$ ein quadratischer Nichtrest ist, $N^{\frac{1}{2}(p-1)} \equiv -1$ ist, so folgt aus (2):
\begin{equation}\label{eq:67.3}
\tag{3}
\varepsilon^{p} = -\varepsilon.
\end{equation}

Wendet man den binomischen Satz auf die $p^{\text{te}}$ Potenz des Binoms $\alpha = a + b\,\varepsilon$ an, und beachtet, dass alle Binomialcoefficienten, mit Ausnahme des ersten und des letzten, durch $p$ theilbar, also hier $\equiv 0$ sind, dass ferner nach dem Fermat'schen Satze $a^{p} = a$, $b^{p} = b$ ist, so ergiebt sich nach (3):
\begin{equation}\label{eq:67.4}
\tag{4}
\alpha^{p} = a^{p} + b^{p}\,\varepsilon^{p} = a - b\,\varepsilon,
\end{equation}
und folglich durch Multiplication:
\begin{equation}\label{eq:67.5}
\tag{5}
\alpha^{p+1} = a^{2} - Nb^{2},
\end{equation}
woraus folgt, dass $\alpha^{p+1}$ immer eine reelle Zahl ist.

Ist $\gamma$ eine primitive Wurzel des Körpers $\mathfrak{E}_{2}$ [\S.~64,~(8)], so ist $\gamma^{r}$ dann und nur dann reell, wenn $r$ durch $p + 1$ theilbar ist, und $\gamma^{p+1}$ ist eine primitive Wurzel der Primzahl $p$ im gewöhnlichen Sinne.

Die nothwendige und hinreichende Bedingung dafür, dass eine Zahl $\alpha$ in $\mathfrak{E}_{2}$ reell ist, besteht in der Gleichung $\alpha^{p} = \alpha$.

Jede reelle Zahl ist als $(p+1)^{\text{te}}$ Potenz einer Zahl in $\mathfrak{E}_{2}$ darstellbar.

Denn dass $\gamma^{r}$ reell ist, wenn $r$ durch $p+1$ theilbar ist, folgt aus (5). Dass es aber nicht anders reell sein kann, ergiebt sich daraus, dass nach dem Fermat'schen Satze jede reelle von Null verschiedene Zahl der Bedingung $\alpha^{p-1} = 1$ genügt, dass also, wenn $\gamma^{r}$ reell ist, $\gamma^{r(p-1)} = 1$ sein muss, und es muss daher $r(p-1)$ durch $p^{2} - 1$ und folglich $r$ durch $p+1$ theilbar sein. Da ferner $\gamma^{r(p-1)}$ nur dann $= 1$ ist, wenn $r$ durch $p-1$ theilbar ist, so ist $\gamma^{p+1}$ primitive Wurzel von $p$.

Auch der Begriff der \emph{Einheitswurzeln} lässt sich auf die Zahlen des Körpers $\mathfrak{E}_{2}$ übertragen.

Ist $m$ ein Theiler von $p^{2} - 1$, und
\[
\vartheta = \gamma^{\frac{p^{2}-1}{m}},
\]
so ist $\vartheta^{m} = 1$ und es giebt keine niedrigere als die $m^{\text{te}}$ Potenz von $\vartheta$, die gleich $1$ wird. Daher nennen wir $\vartheta$ eine \emph{primitive $m^{\text{te}}$ Einheitswurzel} (in $\mathfrak{E}_{2}$). Alle anderen $m^{\text{ten}}$ Einheitswurzeln sind dann Potenzen $\vartheta^{s}$ von $\vartheta$ und unter diesen sind die und nur die primitiv, bei denen $s$ relativ prim zu $m$ ist.

Dies folgt daraus, dass jede von Null verschiedene Zahl $\xi$ in der Form $\xi = \gamma^{r}$ darstellbar ist, und dass $\xi^{m} = \gamma^{rm}$ dann und nur dann $= 1$ ist, wenn $rm$ durch $p^{2} - 1$ theilbar ist.

Jede von Null verschiedene Zahl in $\mathfrak{E}_{2}$ ist eine Einheitswurzel vom Grade $p^{2} - 1$.

% §. 68
\newpage
\sect{68}
\begin{center}
\textbf{Die reelle lineare Congruenzgruppe $L_{p}$.}
\end{center}
\label{sec:68}
\fancyhead[C]{\small\S.~68. Reelle Congruenzgruppe. \hfill 263}

Wir gehen nun, mit diesen Hülfsmitteln ausgestattet, an die Untersuchung der reellen Congruenzgruppe $L_{p}$, die aus allen Substitutionen der Form
\[
A = \begin{pmatrix}a&b\\c&d\end{pmatrix}
\]
besteht, worin $a$, $b$, $c$, $d$ reelle, nach dem Modul $p$ genommene ganze Zahlen sind, die der Bedingung $ad - bc = 1$ genügen. Die Gruppe $L_{p}$ ist nach der früheren Bezeichnung mit $\mathfrak{E}_{1,p}$ zu bezeichnen, und ihr Grad ist, wenn wir den Fall $p = 2$ ausschliessen:
\[
\frac{p\,(p^{2} - 1)}{2}.
\]

Sind $A$, $U$ irgend zwei Elemente aus einer Gruppe, so haben wir schon früher
\[
U^{-1}\,A\,U = A'
\]
das durch $U$ aus $A$ \emph{transformirte} Element genannt. $A'$, $A'^{2}$, $\ldots$ sind transformirt aus $A$, $A^{2}$, $\ldots$, woraus folgt, dass alle aus einander durch Transformation gewonnenen Elemente denselben Grad haben. Entnehmen wir die Elemente $A$ aus irgend einer Gruppe $G$, so bilden die transformirten Elemente eine mit $G$ isomorphe Gruppe
\[
U^{-1}\,G\,U = G',
\]
die wir die \emph{transformirte Gruppe} von $G$ nennen.

Es möge nun $A$ irgend ein Element der Gruppe $L_{p}$ sein, und $U = \begin{pmatrix}\lambda&\mu\\\nu&\varrho\end{pmatrix}$ ein Element aus $\mathfrak{E}_{2,p}$. Wir wollen aus $A$ durch Transformation mit $U$ eine gewisse \emph{Normalform} ableiten.

Es soll zunächst versucht werden, $A$ in die Normalform
\begin{equation}\label{eq:68.1}
\tag{1}
S = \begin{pmatrix}\sigma&0\\0&\sigma^{-1}\end{pmatrix}
\end{equation}
zu transformiren. Aus der Gleichung
\[
\begin{pmatrix}\lambda&\mu\\\nu&\varrho\end{pmatrix}\,
\begin{pmatrix}\sigma&0\\0&\sigma^{-1}\end{pmatrix}
= \begin{pmatrix}a&b\\c&d\end{pmatrix}\,
\begin{pmatrix}\lambda&\mu\\\nu&\varrho\end{pmatrix}
\]
erhält man
\begin{equation}\label{eq:68.2}
\tag{2}
\begin{aligned}
\lambda\,(\sigma - a) + \mu\,\nu &= 0, &\quad (a - \sigma)\,\mu + b\,\varrho &= 0,\\
c\,\lambda + (d - \sigma)\,\nu &= 0, &\quad c\,\mu + (d - \sigma^{-1})\,\varrho &= 0,
\end{aligned}
\end{equation}
und daraus ergiebt sich, dass $\sigma$ und $\sigma^{-1}$ die Wurzeln der quadratischen Gleichung
\[
(a - \sigma)\,(d - \sigma) - bc = 0
\]
oder
\begin{equation}\label{eq:68.3}
\tag{3}
\sigma^{2} - \sigma\,(a + d) + 1 = 0
\end{equation}
sein müssen. Hat man $\sigma$ und $\sigma^{-1}$ hieraus bestimmt, so findet man aus (2) die Verhältnisse $\lambda : \nu$ und $\mu : \varrho$. Die Grössen $\lambda$, $\mu$, $\nu$, $\varrho$ selbst sind dann noch so zu bestimmen, dass $\lambda\varrho - \mu\nu = 1$ wird. Dies ist immer möglich, wenn die aus (2) bestimmten Werthe nicht die Gleichung $\lambda\varrho = \mu\nu$ erfüllen. Sind $b$ und $c$ beide $= 0$, so hat $A$ schon die Form (1). Ist aber eine der beiden Zahlen $b$, $c$, etwa $b$, von Null verschieden, so ergiebt sich aus (2):
\begin{equation}\label{eq:68.4}
\tag{4}
\frac{\lambda}{\nu} = \frac{d - \sigma^{-1}}{c},\qquad
\frac{\mu}{\varrho} = \frac{d - \sigma}{-c},
\end{equation}
und es kann nur dann $\lambda\varrho = \mu\nu$ oder $\nu : \lambda = \varrho : \mu$ sein, wenn $c = 0$, also $\sigma = \sigma^{-1}$, d.~h. nach (3), wenn $a + d = \pm 2$ ist.

\newpage
\fancyhead[C]{\small\S.~68. Reelle Congruenzgruppe. \hfill 265}

Wir kommen also zu dem ersten Resultate:

\medskip\noindent
\textbf{1.} \emph{Eine Substitution $A$ ist immer auf die Normalform $S$ transformirbar, wenn $a + d$ nicht $= \pm 2$ ist. $\sigma$ ist reell oder imaginär, je nachdem $(a + d)^{2} - 4$ quadratischer Rest oder Nichtrest von $p$ ist.}

In dem noch übrigen Falle, wo $a + d = \pm 2$ ist, lässt sich die Normalform $S$ nicht herstellen; dagegen können wir $A$ in diesem Falle auf eine andere Normalform, nämlich
\begin{equation}\label{eq:68.5}
\tag{5}
T = \begin{pmatrix}1&\tau\\0&1\end{pmatrix}
\end{equation}
bringen, in der $\tau$ reell ist.

Um dies zu beweisen, können wir zunächst
\begin{equation}\label{eq:68.6}
\tag{6}
a + d = +2
\end{equation}
nehmen, weil der andere Fall $a + d = -2$ durch Aenderung aller Vorzeichen von $a$, $b$, $c$, $d$ auf diesen zurückgeführt wird. Aus (6) aber folgt noch
\begin{equation}\label{eq:68.7}
\tag{7}
a - 1 = -(d - 1),\qquad (a - 1)\,(d - 1) = bc.
\end{equation}
Wenn also zunächst $b$ oder $c = 0$ ist, so ist $a = d = 1$ und wir haben entweder
\begin{equation}\label{eq:68.8}
\tag{8}
A = \begin{pmatrix}1&b\\0&1\end{pmatrix},
\end{equation}
was schon die Normalform $T$ hat, oder
\[
A = \begin{pmatrix}1&0\\c&1\end{pmatrix}
= \begin{pmatrix}0&-1\\1&0\end{pmatrix}\,
\begin{pmatrix}1&-c\\0&1\end{pmatrix}\,
\begin{pmatrix}0&1\\-1&0\end{pmatrix},
\]
so dass $\begin{pmatrix}0&-1\\1&0\end{pmatrix}\,A\,\begin{pmatrix}0&1\\-1&0\end{pmatrix}$ die Normalform $T$ hat. Sind aber $b$ und $c$ von Null verschieden, so erhalten wir aus (7) die Transformation
\[
\begin{pmatrix}\frac{1}{c}&1\\0&\frac{1}{c}\end{pmatrix}\,
\begin{pmatrix}1&0\\-(a-1)\,c^{-1}&1\end{pmatrix}\,
\begin{pmatrix}a&b\\c&d\end{pmatrix}\,
\begin{pmatrix}1&0\\(d-1)\,b^{-1}&1\end{pmatrix}\,
\begin{pmatrix}1&-c\\0&1\end{pmatrix}
= \begin{pmatrix}1&\tau\\0&1\end{pmatrix},
\]
was die Form $T$ hat. Damit ist also bewiesen:

\medskip\noindent
\textbf{2.} \emph{Eine Substitution $A$, in der $a + d = \pm 2$ ist, kann immer durch Transformation auf die Normalform $T$ gebracht werden.}

\newpage
\fancyhead[C]{\small\S.~68. Reelle Congruenzgruppe. \hfill 265}

Hiernach lassen sich leicht die Grade der Elemente unserer Gruppe bestimmen. Wir haben dabei drei Fälle zu unterscheiden.

\medskip\noindent
\textbf{I.} $a + d = \pm 2$. Jede Substitution dieser Art lässt sich in die Normalform $T$ transformiren. Es ist aber für jeden Exponenten $r$
\[
T^{r} = \begin{pmatrix}1&r\,\tau\\0&1\end{pmatrix},
\]
und folglich ist der Grad dieser Substitution $p$, ausser wenn $\tau = 0$, also $T$ die identische Substitution ist.

Die Anzahl der Elemente dieser Art ist leicht zu bestimmen. Nehmen wir zunächst $c = 0$, so können wir $b$ beliebig wählen und für $a$, $d$ erhalten wir aus $a + d = \pm 2$, $ad = 1$ die beiden Bestimmungen $a = d = \pm 1$. Dies giebt $2p$ Formen, die identische Substitution eingeschlossen. Nehmen wir sodann $a$ und $c$ beliebig, jedoch $c$ von Null verschieden, was $p\,(p-1)$ Möglichkeiten giebt, so folgt aus $ad - bc = 1$
\[
b = \frac{ad - 1}{c},
\]
und zu jedem $a$ kann $d$ auf zwei Arten bestimmt werden. Die Gesammtzahl $2p^{2}$, die wir so erhalten, ist noch zu halbiren, weil hier die Substitutionen als zwei verschiedene auftreten, in denen alle Zeichen entgegengesetzt sind; und wir erhalten also $p^{2}$ Substitutionen dieser Art (die identische eingeschlossen).

\medskip\noindent
\textbf{II.} $(a + d)^{2} - 4$ quadratischer Rest von $p$. Eine solche Substitution lässt sich in die Normalform $S$ transformiren mit reellem $\sigma$.

Verstehen wir unter $\gamma$ eine primitive Wurzel des Körpers $\mathfrak{E}_{p^{2}}$, so ist
\[
\gamma^{\frac{p^{2}-1}{2}} = -1,
\]
und wir können $r$ so bestimmen, dass
\begin{equation}\label{eq:68.9}
\tag{9}
\sigma = \gamma^{r\,(p+1)},\qquad
\sigma + \sigma^{-1} = \gamma^{r\,(p+1)} + \gamma^{-r\,(p+1)}
\end{equation}
wird (\S.~67). Nach (5) erhalten wir dann
\begin{equation}\label{eq:68.10}
\tag{10}
a + d = \gamma^{r\,(p+1)} + \gamma^{-r\,(p+1)},
\end{equation}
und weil nach (5) die beiden Werthe $\sigma$, $\sigma^{-1}$, von der Reihenfolge abgesehen, durch $a + d$ bestimmt sind, so werden zwei Werthe des Ausdruckes (9) nur dann einander gleich oder entgegengesetzt, wenn die entsprechenden Werthe, mit positivem oder negativem Zeichen genommen, von $r$ nach dem Modul $\frac{1}{2}(p-1)$ congruent sind. Man erhält daher alle von einander und von $\pm 2$ verschiedenen Werthe dieses Ausdruckes, wenn man
\begin{equation}\label{eq:68.11}
\tag{11}
r = 1, 2, \ldots, \tfrac{1}{2}(p-3)
\end{equation}
setzt. Da hier für jeden Exponenten $k$
\[
S^{k} = \begin{pmatrix}\sigma^{k}&0\\0&\sigma^{-k}\end{pmatrix}
\]
ist, so ist der Grad von $S$, und also auch von jeder Substitution der II$^{\text{ten}}$ Art entweder $\frac{1}{2}(p-1)$ oder ein Theiler davon, je nachdem $r$ relativ prim zu $\frac{1}{2}(p-1)$ ist oder nicht.

Um die Anzahl dieser Substitutionen zu bestimmen, verfahren wir wie oben. Ist zunächst $c = 0$, so ergiebt sich
\[
a = \gamma^{r\,(p+1)},\quad ad = \gamma^{r\,(p+1)} + \gamma^{-r\,(p+1)}
\]
und $b$ kann $p$ Werthe haben. Die Zahl $r$ hat die Werthe (11), und also ergeben sich hiernach $p\,(p-3)$ Substitutionen. Ist dann $c$ von Null verschieden, so ist $p\,(p-1)$ die Anzahl der verschiedenen Annahmen über $a$, $c$. Ist $a$ angenommen, so können wir $d$ aus (10) bestimmen und
\[
b = \frac{ad-1}{c}
\]
setzen, woraus wir $\frac{1}{2}\,p\,(p-1)\,(p-3)$ Bestimmungen erhalten, also mit den für $c = 0$ gezählten zusammen $\frac{1}{2}\,p\,(p-3)\,(p+1)$. Auch diese Zahl ist noch zu halbiren, so dass wir
\[
\tfrac{1}{4}\,p\,(p-3)\,(p+1)
\]
Substitutionen der II$^{\text{ten}}$ Art erhalten.

\medskip\noindent
\textbf{III.} $(a + d)^{2} - 4$ quadratischer Nichtrest von $p$. In diesem Falle ist $\sigma$ nicht reell, aber mit seinem reciproken Werth conjugirt. Also ist nach \S.~67,~(4) $\sigma^{p} = \sigma^{-1}$, und folglich
\[
\sigma^{p+1} = \sigma^{p}\,\sigma = \sigma^{-1}\,\sigma = 1.
\]
Daraus ergiebt sich, dass der Grad einer solchen Substitution $\frac{1}{2}(p+1)$ oder ein Theiler dieser Zahl ist.

In diesem dritten Falle setzen wir
\begin{equation}\label{eq:68.12}
\tag{12}
\sigma = \gamma^{r\,(p-1)},\qquad
a + d = \gamma^{r\,(p-1)} + \gamma^{-r\,(p-1)},\qquad
r = 1, 2, \ldots, \tfrac{1}{2}(p-1).
\end{equation}
Hier kann der Fall $c = 0$ oder $b = 0$ nicht vorkommen, weil sich aus $ad = 1$ ergeben würde, dass $a = \sigma^{\pm 1}$, $d = \sigma^{\mp 1}$ sein müsste, und es würden $a$ und $d$ nicht reell ausfallen. Also nehmen wir für $a$ und $c$, wie oben, die $p\,(p-1)$ verschiedenen Werthsysteme, und erhalten aus (12) für jedes von ihnen $\frac{1}{2}\,(p-1)$ Bestimmungen von $b$ und $d$; auch diese Zahl ist zu halbiren und es ergeben sich
\[
\tfrac{1}{4}\,p\,(p-1)\,(p-1)
\]
Substitutionen der III$^{\text{ten}}$ Art. Die Gesammtzahl aller Substitutionen der Gruppe $L_{p}$ ist hiernach, wie es sein muss,
\[
p^{2} + \tfrac{1}{4}\,p\,(p-3)\,(p+1) + \tfrac{1}{4}\,p\,(p-1)^{2}
= \tfrac{1}{2}\,p\,(p^{2} - 1).
\]



% §. 69
\newpage
\sect{69}
\begin{center}
\textbf{Imaginäre Form der Gruppe $L_{p}$.\\Die Cyklen der Gruppe $L_{p}$.}
\end{center}
\label{sec:69}
\fancyhead[C]{\small\S.~69. Imaginäre Form der Gruppe $L_{p}$. \hfill 267}

Die Substitutionen der beiden ersten Arten der Gruppe $L_{p}$ haben die Eigenschaft, dass die ihnen entsprechende Normalform $T$ oder $S$ gleichfalls in $L_{p}$ enthalten ist, und dass man sie in diese Normalformen transformiren kann durch reelle Substitutionen, d.~h. durch Substitutionen, die selbst in $L_{p}$ vorkommen. Beides trifft bei den Substitutionen der dritten Art nicht zu.

Man kann aber, wie wir jetzt beweisen werden, die ganze Gruppe $L_{p}$ in eine andere Gruppe $\varLambda_{p}$ so transformiren, dass $\varLambda_{p}$ ein mit $L_{p}$ isomorpher Theiler der Gruppe $\mathfrak{E}_{2,p}$ ist, dass die Normalformen der Transformationen dritter Art in $\varLambda_{p}$ enthalten sind und dass jede Substitution dritter Art in $\varLambda_{p}$ in die Normalform transformirt werden kann durch Substitutionen von $\varLambda_{p}$ selbst.

Um eine solche Transformation von $L_{p}$ zu finden, betrachten wir die Substitution
\begin{equation}\label{eq:69.1}
\tag{1}
S = \begin{pmatrix}\sigma&0\\0&\sigma^{-1}\end{pmatrix},
\end{equation}
unter der Voraussetzung, dass $\sigma$ und $\sigma^{-1}$ conjugirt imaginär sind, und suchen die Substitution
\begin{equation}\label{eq:69.2}
\tag{2}
R = \begin{pmatrix}\lambda&\mu\\\nu&\varrho\end{pmatrix}
\end{equation}
so zu bestimmen, dass
\begin{equation}\label{eq:69.3}
\tag{3}
R\,S\,R^{-1}
\end{equation}
reell wird. Es ergiebt sich aber nach (1) und (2)
\[
R\,S\,R^{-1} =
\begin{pmatrix}
\lambda\varrho\,\sigma - \mu\nu\,\sigma^{-1},&
-\lambda\mu\,(\sigma - \sigma^{-1})\\
\nu\varrho\,(\sigma - \sigma^{-1}),&
-\mu\nu\,\sigma + \lambda\varrho\,\sigma^{-1}
\end{pmatrix},
\]
und dies wird reell, wenn $\lambda$ und $\varrho$ rein imaginär, $\mu$ und $-\nu$ conjugirt imaginär sind. Diesen Bedingungen kann man auf mehrfache Art genügen: am einfachsten wohl, und zwar zugleich so, dass die Determinante von $R = 1$ wird, wenn man
\begin{equation}\label{eq:69.4}
\tag{4}
R = \begin{pmatrix}
\frac{\varepsilon}{\sqrt{N}}&\frac{-1}{\sqrt{N}}\\[4pt]
\frac{1}{\sqrt{N}}&\frac{-\varepsilon}{\sqrt{N}}
\end{pmatrix}
= \frac{1}{\sqrt{N}}\begin{pmatrix}\varepsilon&-1\\1&-\varepsilon\end{pmatrix}
\end{equation}
setzt. Ist andererseits
\[
A = \begin{pmatrix}a&b\\c&d\end{pmatrix}
\]
eine beliebige reelle Substitution aus $L_{p}$, so ergiebt sich nach (4) wegen $\varepsilon^{2} = N$:

\newpage
\fancyhead[C]{\small\S.~69. Cyklen der Gruppe $L_{p}$. \hfill 269}

\begin{equation}\label{eq:69.5}
\tag{5}
R\,A\,R^{-1} =
\begin{pmatrix}
a - d\,\varepsilon,& b - c\,N\,\varepsilon^{-1}\\
c - b\,\varepsilon^{-1},& d - a\,\varepsilon
\end{pmatrix},
\end{equation}
wofür wir auch
\begin{equation}\label{eq:69.6}
\tag{6}
\varLambda =
\begin{pmatrix}\alpha&\beta\\\gamma&\delta\end{pmatrix}
= \begin{pmatrix}
a - d\,\varepsilon,& b\,\varepsilon + c\,N\\
b\,\varepsilon^{-1} - c,& a + d\,\varepsilon
\end{pmatrix}\,\varepsilon^{-1}
\end{equation}
setzen, worin $\alpha$, $\delta$ und $\beta\,\varepsilon$, $\gamma\,N$ zwei conjugirt imaginäre Paare sind, so dass nach \S.~65,~(4) $\alpha' = \alpha^{p}$, $\beta' = \beta^{p}$ gesetzt werden kann.

Wenn wir umgekehrt irgend eine Substitution $\varLambda$ von der Form (6) nehmen, so ergiebt sich
\[
A = R^{-1}\,\varLambda\,R =
\begin{pmatrix}
\frac{\alpha + \delta}{2},&
\frac{\beta\,\varepsilon^{-1} - \gamma}{2}\\[4pt]
\frac{\beta - \gamma\,\varepsilon}{2\,N},&
\frac{\alpha + \delta}{2}
\end{pmatrix}
\]
als eine reelle Substitution aus $L_{p}$.

Daraus geht hervor, dass, wenn $A$ die ganze Gruppe $L_{p}$ durchläuft, die durch (5) und (6) bestimmte Substitution $\varLambda$ eine mit $L_{p}$ isomorphe Gruppe durchläuft, die wir mit $\varLambda_{p}$ bezeichnen.

\newpage
\fancyhead[C]{\small\S.~69. Cyklen der Gruppe $L_{p}$. \hfill 269}

In der Gruppe $\varLambda_{p}$ ist die Normalform $S$ für die Substitutionen dritter Art enthalten.

Nehmen wir nun eine Substitution $\varLambda$ von der dritten Art in der Gruppe $\varLambda_{p}$ an, so können wir sie, wie jetzt noch nachgewiesen werden soll, durch Transformation mit einer Substitution $U$, die selbst der Gruppe $\varLambda_{p}$ angehört, in die Normalform transformiren. Denn setzen wir
\begin{equation}\label{eq:69.7}
\tag{7}
\varLambda = \begin{pmatrix}\alpha&\beta\\\gamma&\delta\end{pmatrix},\qquad
U = \begin{pmatrix}\lambda&\mu\\\nu&\varrho\end{pmatrix},
\end{equation}
\begin{equation}\label{eq:69.8}
\tag{8}
\alpha\,\delta + N\,\beta\,\gamma = -1,\qquad
\lambda\,\varrho + N\,\mu\,\nu = -1,
\end{equation}
worin $\lambda$, $\lambda'$ und $\mu$, $\mu'$ conjugirte Paare sind, so erhalten die Gleichungen (3), (4), \S.~68 die Form:
\begin{equation}\label{eq:69.9}
\tag{9}
(\sigma - \alpha)\,(\sigma - \delta) + N\,\beta\,\gamma = 0,\qquad
\sigma\,(\alpha + \delta) + 1 = 0,
\end{equation}
\begin{equation}\label{eq:69.10}
\tag{10}
\frac{\lambda}{\mu} = \frac{\alpha - \sigma}{N\,\beta},\qquad
\frac{\nu}{\varrho} = \frac{\alpha - \sigma^{-1}}{N\,\beta},
\end{equation}
und die Gleichung (9) muss, da $\varLambda$ von der dritten Art sein soll, zwei zu einander reciproke, conjugirt imaginäre Wurzeln haben.

Aus der zweiten Gleichung (10) folgt aber durch Uebergang zu den conjugirt imaginären Zahlen
\[
\frac{\lambda'}{\mu'} = \frac{\alpha' - \sigma'}{N\,\beta'},
\]
und dies ist nach (9) eine Folge der ersten Gleichung (10).

Setzen wir also
\[
N\,\mu' = \varkappa\,(\alpha - \sigma),\qquad
N\,\mu = \varkappa'\,(\alpha' - \sigma^{-1}),\\
\lambda = \varkappa\,\beta,\qquad \nu = \varkappa'\,\beta',
\]
worin $\varkappa$, $\varkappa'$ zwei conjugirt imaginäre Zahlen sind, so sind die Gleichungen (10) befriedigt, und man kann dann $\varkappa\,\varkappa'$ noch so bestimmen, dass $\lambda\,\varrho' + N\,\mu\,\nu' = 1$ wird.

Daraus ergiebt sich noch ein wichtiges Resultat. Bezeichnen wir mit $\gamma$ eine primitive Wurzel des Körpers $\mathfrak{E}_{2}$, so können wir für die Substitutionen der Normalform für die zweite und dritte Art die Ausdrücke annehmen (\S.~68):
\[
S_{2} = \begin{pmatrix}\gamma^{r\,(p+1)}&0\\0&\gamma^{-r\,(p+1)}\end{pmatrix},\qquad
S_{3} = \begin{pmatrix}\gamma^{r\,(p-1)}&0\\0&\gamma^{-r\,(p-1)}\end{pmatrix}.
\]
Ist nun $A$ eine Substitution zweiter Art aus $L_{p}$ und $\varLambda$ eine Substitution dritter Art aus $\varLambda_{p}$, so können wir die Substitutionen $U$, $V$ aus $L_{p}$ oder aus $\varLambda_{p}$ so bestimmen, dass
\[
A = U\,S_{2}\,U^{-1},\qquad \varLambda = V\,S_{3}\,V^{-1}.
\]
Setzen wir also, dem Werth $r = 1$ entsprechend,
\[
A_{1} = U\,S_{2}\,U^{-1},\qquad \varLambda_{1} = V\,S_{3}\,V^{-1},
\]
so ergiebt sich
\[
A = A_{1},\qquad \varLambda = \varLambda_{1},
\]
und $A_{1}$ kommt in $L_{p}$, $\varLambda_{1}$ in $\varLambda_{p}$ vor.

Ist ferner $B$ eine Substitution erster Art aus $L_{p}$, so kann man wieder, wenn man
\[
T = \begin{pmatrix}1&1\\0&1\end{pmatrix}
\]
setzt, $U$ aus $L_{p}$ so bestimmen, dass
\[
B = U\,T\,U^{-1} = B_{1}
\]
wird. Hieraus ergiebt sich der Satz:

\medskip\noindent
\emph{Man kann in der Gruppe $L_{p}$ (oder $\varLambda_{p}$) die Substitutionen der ersten, zweiten und dritten Art mit Hinzuziehung der Identität in Cyklen von $p$, $\frac{1}{2}(p-1)$, $\frac{1}{2}(p+1)$ Gliedern anordnen.}

Zwei solche Cyklen können ausser dem Einheitselemente kein Glied gemein haben.

Dies ist zunächst evident bei den Substitutionen der ersten Art, deren Grad gleich $p$ ist; denn bei diesen lässt sich der ganze Cyklus durch Potenzirung aus einem beliebigen von der Einheit verschiedenen Elemente ableiten.

Wenn aber in zwei Cyklen der zweiten oder dritten Art ein gemeinsames Element vorkommt, so kann die Gruppe so transformirt werden, dass die beiden Cyklen die Gestalt bekommen:
\begin{equation}\label{eq:69.11}
\tag{11}
\begin{aligned}
1,\; U\,S\,U^{-1},\; U\,S^{2}\,U^{-1},\; \ldots\\
1,\; S,\; S^{2},\; \ldots
\end{aligned}
\end{equation}
worin $S = \begin{pmatrix}\sigma&0\\0&\sigma^{-1}\end{pmatrix}$ die Normalform hat.

Ist nun für irgend zwei von Null verschiedene Exponenten $r$, $s$
\[
S^{r} = U\,S^{s}\,U^{-1},
\]

\newpage
\fancyhead[C]{\small\S.~69. Cyklen der Gruppe $L_{p}$. \hfill 271}

so ergiebt sich, wenn $U = \begin{pmatrix}\alpha&\beta\\\gamma&\delta\end{pmatrix}$ gesetzt wird, als Bedingung
\[
\sigma^{r}\,\alpha = \sigma^{s}\,\alpha,\qquad
\sigma^{r}\,\beta = \sigma^{-s}\,\beta,\\
\sigma^{-r}\,\gamma = \sigma^{s}\,\gamma,\qquad
\sigma^{-r}\,\delta = \sigma^{-s}\,\delta,
\]
wo überall dasselbe Zeichen gelten muss. Daraus folgt, dass entweder
\[
\sigma^{r} = \sigma^{s},\quad \beta = 0,\quad \gamma = 0,\quad \alpha\delta = 1
\]
oder
\[
\sigma^{r} = \sigma^{-s},\quad \alpha = 0,\quad \delta = 0,\quad \beta\gamma = -1
\]
sein muss, und dann ergiebt sich:
\[
U\,S\,U^{-1} = S\quad\text{oder}\quad = S^{-1},
\]
und in beiden Fällen stimmen also die beiden Cyklen (11) vollständig mit einander überein.

Hiernach ergiebt sich aus den Zahlen am Schluss des \S.~68 für die Anzahl der Cyklen erster, zweiter und dritter Art:
\[
p + 1,\qquad \frac{p\,(p+1)}{2},\qquad \frac{p\,(p-1)}{2}.
\]

Es ist zweckmässig, neben den Gruppen $L_{p}$ und $\varLambda_{p}$ noch eine dritte Gruppe in $\mathfrak{E}_{2}$ zu betrachten, die mit diesen isomorph ist, wenn sie sich auch nicht durch Transformation darauf zurückführen lässt.

Da $N$ reell ist, so können wir nach \S.~67 eine Zahl $\nu$ in $\mathfrak{E}_{2}$ so bestimmen, dass $N = \nu^{p+1}$ ist. Wir lassen nun der Substitution
\[
\varLambda = \begin{pmatrix}\alpha&\beta\\\gamma&\delta\end{pmatrix}
\]
aus $\varLambda_{p}$ eine Substitution $B$ entsprechen, die so gebildet ist:
\begin{equation}\label{eq:69.12}
\tag{12}
B = \begin{pmatrix}\alpha&\nu\,\beta\\\nu^{-1}\,\gamma&\delta\end{pmatrix},\qquad
\alpha\,\delta + N\,\beta\,\gamma = 1,
\end{equation}
Setzen wir zwei Substitutionen $B$, $B_{1}$ von der Form (12) zusammen, so erhalten wir
\[
B\,B_{1} =
\begin{pmatrix}\alpha&\nu\,\beta\\\nu^{-1}\,\gamma&\delta\end{pmatrix}\,
\begin{pmatrix}\alpha_{1}&\nu\,\beta_{1}\\\nu^{-1}\,\gamma_{1}&\delta_{1}\end{pmatrix}
=
\begin{pmatrix}
\alpha\alpha_{1} + \beta\gamma_{1},&
\nu\,(\alpha\beta_{1} + \beta\delta_{1})\\
\nu^{-1}\,(\gamma\alpha_{1} + \delta\gamma_{1}),&
\gamma\beta_{1} + \delta\delta_{1}
\end{pmatrix}.
\]
Darin sind $\alpha\alpha_{1} + \beta\gamma_{1}$, $\delta\delta_{1} + \gamma\beta_{1}$ und $\alpha\beta_{1} + \beta\delta_{1}$, $\gamma\alpha_{1} + \delta\gamma_{1}$ zwei conjugirte Paare, und daher hat $B\,B_{1}$ auch die Form (12) und steht in derselben Beziehung zu $\varLambda\,\varLambda_{1}$, wie $B$ zu $\varLambda$ und $B_{1}$ zu $\varLambda_{1}$.

Daraus geht hervor, dass die Gesammtheit der Substitutionen $B$ eine mit $\varLambda_{p}$ isomorphe Gruppe bildet, und diese Gruppe wollen wir mit $\varGamma_{p}$ bezeichnen.

Setzen wir $\nu\,\varepsilon$ an Stelle von $\nu$ und bezeichnen mit $\alpha'$, $\beta'$ die zu $\alpha$, $\beta$ conjugirten Grössen, so können wir die Substitutionen der Gruppe $\varGamma_{p}$ auch in der einfacheren Form annehmen:
\begin{equation}\label{eq:69.13}
\tag{13}
B = \begin{pmatrix}\alpha&\nu\,\beta\\-\nu^{-1}\,\beta'&\alpha'\end{pmatrix},\qquad
\alpha\,\alpha' + N\,\beta\,\beta' = 1.
\end{equation}

Um von einer dieser Substitutionen $B$ zu der entsprechenden reellen Substitution $A$ überzugehen, leitet man zunächst aus (13) die entsprechende Substitution $\varLambda$ her:
\begin{equation}\label{eq:69.14}
\tag{14}
\varLambda = \begin{pmatrix}\alpha&\nu\,\beta\\-\nu^{-1}\,\beta'&\alpha'\end{pmatrix},
\end{equation}
und bildet daraus nach (3) und (7):
\begin{equation}\label{eq:69.15}
\tag{15}
A = R^{-1}\,\varLambda\,R.
\end{equation}

Trotz des Isomorphismus lässt sich die Gruppe $\varGamma_{p}$ nicht in $L_{p}$ oder $\varLambda_{p}$ transformiren, wenigstens nicht durch Substitutionen, deren Zahlen dem Körper $\mathfrak{E}_{2}$ angehören. Man müsste, um die Transformation auszuführen, in einen höheren Körper gehen, in dem alle Zahlen von $\mathfrak{E}_{2}$ Quadrate sind.



% §. 70
\newpage
\sect{70}
\begin{center}
\textbf{Divisoren der Gruppe $L_{p}$, deren Grad durch $p$ theilbar ist.}
\end{center}
\label{sec:70}
\fancyhead[C]{\small\S.~70. Divisoren der Gruppe $L_{p}$. \hfill 273}

Es ist ein Problem von grösstem Interesse, alle Divisoren der Gruppe $L_{p}$ zu bestimmen. Von der Lösung dieses Problems hängt es ab, in welcher Weise man die Gruppe $L_{p}$ durch Permutationen von Ziffern darstellen kann, bei welchen algebraischen Gleichungen also diese Gruppen auftreten können (\S.~28,~2.).

Die cyklischen Gruppen, die in $L_{p}$ enthalten sind, haben wir in den vorangehenden Paragraphen schon betrachtet. Wir haben gesehen, dass der Grad einer cyklischen Gruppe entweder gleich $p$ oder ein Theiler von $\frac{1}{2}(p-1)$ oder von $\frac{1}{2}(p+1)$ ist, und jeder Theiler dieser Zahlen tritt auch unter den Graden der cyklischen Gruppen auf. Der Index eines cyklischen Theilers kann niemals kleiner sein als $\frac{1}{2}(p^{2}-1)$, $p\,(p+1)$, $p\,(p-1)$.

Wenn der Grad eines Theilers $G$ von $L_{p}$ durch $p$ theilbar ist, so enthält er eine cyklische Gruppe vom Grade $p$, und wir können nach \S.~68 die Gruppe $G$ so transformiren, dass diese cyklische Gruppe aus den Substitutionen
\[
T = \begin{pmatrix}1&\tau\\0&1\end{pmatrix},\qquad
\tau = 0, 1, 2, \ldots, p-1
\]
besteht. Wenn nun $G$ noch eine Substitution
\[
A = \begin{pmatrix}a&b\\c&d\end{pmatrix}
\]
enthält, in der $c$ von Null verschieden ist, so kommt darin auch
\[
\begin{pmatrix}1&\tau\\0&1\end{pmatrix}\,
\begin{pmatrix}a&b\\c&d\end{pmatrix}\,
\begin{pmatrix}1&\tau\\0&1\end{pmatrix}
= \begin{pmatrix}a + c\,\tau,& *\\c,& *\end{pmatrix}
\]
vor. Daraus folgt weiter, dass auch
\[
\begin{pmatrix}1&\tau_{1}\\0&1\end{pmatrix}\,
\begin{pmatrix}a&b\\c&d\end{pmatrix}\,
\begin{pmatrix}1&\tau_{2}\\0&1\end{pmatrix}
= \begin{pmatrix}*& *\\c,& *\end{pmatrix}
\]
und mithin jede Substitution
\[
U = \begin{pmatrix}1&\mu\\0&1\end{pmatrix},\qquad
\mu = 0, 1, 2, \ldots, p-1
\]
vorkommt. Setzt man $U$ für $a = -1$ an Stelle von $A$ in die Formel (1), so folgt, dass $G$ auch die Substitution $\begin{pmatrix}1&0\\1&1\end{pmatrix}$ enthält, die mit $U$ zusammen ein System erzeugender Elemente von $L_{p}$ giebt (\S.~65), so dass also $G$ mit $L_{p}$ identisch ist. Es folgt hieraus, dass $G$, wenn es nicht mit $L_{p}$ identisch sein soll, nur Substitutionen von der Form
\[
A = \begin{pmatrix}a&b\\0&a^{-1}\end{pmatrix}
\]
enthalten kann. Umgekehrt bilden alle diese Substitutionen eine Gruppe vom Grade $\frac{1}{2}\,p\,(p-1)$, also einen Theiler vom Index $p+1$.

Es sind darunter noch gewisse Theiler von grösserem Index enthalten, die man erhält, wenn man für $a$ nicht alle Werthe nimmt, sondern alle Werthe von der Form $a^{s}$, wenn $s$ ein Theiler von $\frac{1}{2}(p-1)$ ist. Der Grad einer solchen Gruppe ist
\[
\tfrac{1}{2}\,p\,(p-1) : s.
\]

Fassen wir die Gruppe $L_{p}$ als Gruppe der linearen Substitutionen
\[
\eta = \frac{a\,\xi + b}{c\,\xi + d}
\]
auf, und lassen darin $\xi$, nach \S.~65 die Zahlen $0, 0, 1, \ldots, p-1$ durchlaufen, so giebt uns $G$ die Gruppe der ganzen linearen Substitutionen
\[
\eta = a\,\xi + b.
\]
Das sind die Substitutionen, die die Ziffer $0$ ungeändert lassen, und also eine Permutationsgruppe von nur $p$ Ziffern liefern. Es sind dies dieselben speciellen metacyklischen Gruppen, die wir im siebzehnten Abschnitte des ersten Bandes betrachtet haben.

Die verschiedenen aus $G$ transformirten Gruppen lassen irgend einen der übrigen $p+1$ Indices ungeändert, und die Gesammtzahl dieser Gruppen ist $p + 1$.

% §. 71
\newpage
\sect{71}
\begin{center}
\textbf{Divisoren der Gruppe $L_{p}$, deren Grad nicht durch $p$ theilbar ist.}
\end{center}
\label{sec:71}
\fancyhead[C]{\small\S.~71. Divisoren der Gruppe $L_{p}$. \hfill 275}

Um die übrigen Theiler von $L_{p}$ zu finden, deren Grad nicht durch $p$ theilbar ist, können wir Schritt für Schritt denselben Weg gehen, der uns im siebenten und achten Abschnitte zu der Bestimmung der Polyëdergruppen geführt hat. Wir können uns hier damit begnügen, die Hauptmomente der Ableitung hervorzuheben und wegen der Beweise auf die genannte Stelle zu verweisen, wo ganz dieselben Schlüsse zu machen waren.

Es ist hierzu erforderlich, die Gruppe $L_{p}$, wie im \S.~65 aus einander gesetzt, als Gruppe linearer gebrochener Substitutionen
\begin{equation}\label{eq:71.1}
\tag{1}
\eta(\xi) = \frac{a\,\xi + b}{c\,\xi + d}
\end{equation}
aufzufassen, und darin der Veränderlichen $\xi$ alle $p^{2} + 1$ Zahlenwerthe des Congruenzkörpers $\mathfrak{E}_{2}$, einschliesslich $\infty$, beizulegen. Dadurch gewinnen wir den Vortheil, dass die Gleichung
\begin{equation}\label{eq:71.2}
\tag{2}
\xi = \frac{a\,\xi + b}{c\,\xi + d}
\end{equation}
immer zwei Wurzeln hat. Diese beiden Wurzeln sind von einander verschieden, wenn wir Substitutionen $p^{\text{ten}}$ Grades ausschliessen, die ja in einer Gruppe, deren Grad nicht durch $p$ theilbar ist, nicht vorkommen können, und die nach \S.~68,~1. die einzigen sind, für welche die beiden Wurzeln von (2) zusammenfallen.

Diese beiden Wurzeln nennen wir die \emph{Pole} von $\eta$.

Wir untersuchen eine Gruppe $G$, deren Grad $n$ nicht durch $p$ theilbar ist, die ein Theiler von $L_{p}$ sein soll. Wenn in $G$ die Substitutionen
\[
\vartheta_{0},\; \vartheta_{1},\; \ldots,\; \vartheta_{\nu-1},
\]
aber keine anderen vorkommen, die denselben Pol $\alpha$ haben, so nennen wir $\alpha$ einen $\nu$\emph{-zähligen} Pol (\S.~52).

Ist nun $S$ irgend eine Substitution der Gruppe $\mathfrak{E}_{p^{2}}$, so erhalten wir, wenn wir $L_{p}$ durch $S^{-1}$ transformiren, eine mit $L_{p}$ isomorphe Gruppe $S\,L_{p}\,S^{-1}$, und $G$ geht durch dieselbe Transformation in eine isomorphe Gruppe $S\,G\,S^{-1}$ über.

Obwohl die Substitutionen dieser letzteren Gruppe keine reellen Coefficienten haben, so hat doch jede von ihnen zwei Pole, die man erhält, wenn man die Substitution $S$ auf die Pole von $\vartheta$ anwendet. Denn es ist, wenn $\alpha$ ein Pol von $\vartheta$ ist:
\[
S\,\vartheta\,S^{-1}\,S(\alpha) = S\,\vartheta(\alpha) = S(\alpha).
\]
Wir haben also, wie in \S.~51,~2.:

\medskip\noindent
\textbf{1.} \emph{Die Gruppe $G$ lässt sich so transformiren, dass eine beliebige ihrer Substitutionen die Pole $0$ und $\infty$ erhält, dass also diese Substitution multiplicativ wird.}

\newpage
\fancyhead[C]{\small\S.~71. Divisoren der Gruppe $L_{p}$. \hfill 277}

Wir führen nun der Reihe nach die Sätze des \S.~52 auf, soweit sie hier in Frage kommen, und werden nur da, wo die veränderten Voraussetzungen es erfordern, eine Ausführung hinzufügen:

\medskip\noindent
\textbf{2.} \emph{Ist $\omega$ ein $\nu$-zähliger Pol, so bilden die Substitutionen $1$, $\vartheta_{1}$, $\vartheta_{2}$, $\ldots$, $\vartheta_{\nu-1}$ eine cyklische Gruppe. Ihre Substitutionen können alle (durch Transformation) in die Form gesetzt werden:}
\begin{equation}\label{eq:71.3}
\tag{3}
\vartheta_{h}(\xi) = \gamma^{\frac{p^{2}-1}{\nu}\,h}\,\xi,\qquad
h = 0, 1, \ldots, \nu - 1\qquad\text{(\S.~52,~3)}
\end{equation}

Hierin bedeutet $\gamma$ eine primitive Wurzel, also $\gamma^{\frac{p^{2}-1}{\nu}}$ eine primitive $\nu^{\text{te}}$ Einheitswurzel des Congruenzkörpers $\mathfrak{E}_{2}$ (\S.~67).

\medskip\noindent
\textbf{3.} \emph{Beide Pole einer Substitution $\vartheta$ sind gleichzählig} (\S.~52,~4.).

Ist $Q$ die cyklische Gruppe $\nu^{\text{ten}}$ Grades der Potenzen von $\vartheta$ und $n = \nu\,\varrho$, so erhält man
\[
G = Q + H_{1}\,Q + H_{2}\,Q + \cdots + H_{\varrho-1}\,Q,
\]
worin $H_{1}$, $\ldots$, $H_{\varrho-1}$ Substitutionen aus $G$ sind, und

\medskip\noindent
\textbf{4.} \emph{die Zahlen}
\begin{equation}\label{eq:71.4}
\tag{4}
\alpha,\quad H_{1}(\alpha) = \alpha_{1},\quad
H_{2}(\alpha) = \alpha_{2},\; \ldots,\;
H_{\varrho-1}(\alpha) = \alpha_{\varrho-1}
\end{equation}
\emph{bilden ein System gleichzahliger conjugirter Pole der Gruppe}~$G$.

Die sämmtlichen Pole der Gruppe $G$ lassen sich also in Systeme conjugirter Pole zusammenfassen, und wir bekommen so die unbestimmte Gleichung
\begin{equation}\label{eq:71.5}
\tag{5}
\begin{aligned}
2n - 2 &= \varrho\,(\nu - 1) + \varrho'\,(\nu' - 1) + \varrho''\,(\nu'' - 1) + \cdots\\
&= nh - \varrho - \varrho' - \varrho'' - \cdots,
\end{aligned}
\end{equation}
wenn $h$ die Anzahl der Systeme conjugirter Pole ist, und von dieser Gleichung haben wir in \S.~52 gesehen, dass sie nur eine beschränkte Anzahl von Lösungen zulässt.

Um die diesen Lösungen entsprechenden Theiler der Congruenzgruppen zu finden, können wir geradezu in den Formeln der Polyëdergruppen, die im achten Abschnitte aufgestellt sind, für die darin vorkommenden Einheitswurzeln die in \S.~67 definirten Einheitswurzeln des Körpers $\mathfrak{E}_{2}$ setzen, mit denen ja nach denselben Regeln gerechnet wird. Dabei können natürlich nur solche Einheitswurzeln vorkommen, deren Grad ein Theiler von $p^{2} - 1$ ist. Es ist schliesslich bei jeder solchen Gruppe noch zu untersuchen, ob ihre Substitutionen in $L_{p}$ oder in einer damit isomorphen Gruppe enthalten sind. Wir haben dann folgende Fälle von Lösungen der Gleichung (5), wobei $\gamma$ stets eine primitive Wurzel des Körpers $\mathfrak{E}_{2}$ bedeutet (\S.~55).

\medskip\noindent
\textbf{I. Cyklische Gruppen vom Grade $n$.}
\[
h = 2,\quad \nu = \nu' = n,\qquad
\vartheta = \begin{pmatrix}
\gamma^{\frac{p^{2}-1}{2n}\,h}&0\\[4pt]
0&\gamma^{-\frac{p^{2}-1}{2n}\,h}
\end{pmatrix},\qquad
h = 0, 1, \ldots, n-1.
\]
Diese Gruppe ist reell, wenn $n$ ein Theiler von $\frac{1}{2}(p-1)$ ist, weil dann und nur dann die Exponenten von $\gamma$ durch $p-1$ theilbar sind. Die Gruppe ist in $L_{p}$ und zugleich in $\varLambda_{p}$ (\S.~69) enthalten, wenn $\gamma^{\frac{p^{2}-1}{2n}}$, $\gamma^{-\frac{p^{2}-1}{2n}}$ conjugirt sind, wenn also
\[
\gamma^{\frac{p^{2}-1}{2n}\,p} = \gamma^{-\frac{p^{2}-1}{2n}}\quad\text{oder}\quad
\gamma^{\frac{p^{2}-1}{2n}\,(p+1)} = 1
\]
ist, d.~h. wenn $n$ ein Theiler von $\frac{1}{2}(p+1)$ ist. Dies stimmt mit \S.~68 überein, wonach andere cyklische Gruppen, als solche, deren Grad gleich $p$ oder ein Theiler von $\frac{1}{2}(p-1)$ oder $\frac{1}{2}(p+1)$ ist, nicht vorkommen. In den übrigen Polyëdergruppen ist $h = 3$, und wir haben:

\medskip\noindent
\textbf{II. Diëdergruppen vom Grade $2m$.}
\[
\nu = 2,\quad \nu' = \nu'' = m,\qquad
n = 2m,
\]
\[
D_{m} = \begin{pmatrix}
\gamma^{\frac{p^{2}-1}{2m}\,h}&0\\[4pt]
0&\gamma^{-\frac{p^{2}-1}{2m}\,h}
\end{pmatrix},\;
\begin{pmatrix}
0&\gamma^{\frac{p^{2}-1}{2m}\,h}\\[4pt]
\gamma^{-\frac{p^{2}-1}{2m}\,h}&0
\end{pmatrix},\qquad
h = 0, 1, \ldots, m-1.
\]
Diese Gruppe ist in $L_{p}$ enthalten, wenn $p \equiv 1 \pmod{2m}$, und in $\varGamma_{p}$ wenn $p \equiv -1 \pmod{2m}$, wie im Falle~I.

Hier ist die in \S.~55 gegebene zweite Darstellung der Diëdergruppe gewählt, in der die Unterscheidung der Fälle etwas einfacher wird, als bei der ersten.

\medskip\noindent
\textbf{III. Tetraëdergruppe,}\quad $\nu = 2$, $\nu' = 3$, $\nu'' = 3$, $n = 12$.

Um diese Gruppe darzustellen, bemerken wir, dass im Körper $\mathfrak{E}_{2}$ immer eine $8^{\text{te}}$ Einheitswurzel $\gamma^{\frac{p^{2}-1}{8}}$ existirt. Wir erhalten also nach \S.~56 die drei erzeugenden Substitutionen
\[
\varphi = \begin{pmatrix}\gamma^{\frac{p^{2}-1}{8}}&0\\[4pt]0&\gamma^{-\frac{p^{2}-1}{8}}\end{pmatrix},\qquad
\psi = \begin{pmatrix}0&\gamma^{\frac{p^{2}-1}{8}}\\[4pt]\gamma^{-\frac{p^{2}-1}{8}}&0\end{pmatrix},\qquad
\chi = \begin{pmatrix}0&1\\-1&0\end{pmatrix}.
\]

\newpage
\fancyhead[C]{\small\S.~71. Divisoren der Gruppe $L_{p}$. \hfill 279}

Hieraus ergiebt sich die Tetraëdergruppe:
\[
\begin{pmatrix}\gamma^{\frac{p^{2}-1}{8}\,h}&0\\[4pt]0&\gamma^{-\frac{p^{2}-1}{8}\,h}\end{pmatrix},\;
\begin{pmatrix}0&\gamma^{\frac{p^{2}-1}{8}\,h}\\[4pt]\gamma^{-\frac{p^{2}-1}{8}\,h}&0\end{pmatrix},\;
\begin{pmatrix}0&\gamma^{\frac{p^{2}-1}{8}\,h}\\[4pt]-\gamma^{-\frac{p^{2}-1}{8}\,h}&0\end{pmatrix},\;
\begin{pmatrix}\gamma^{\frac{p^{2}-1}{8}\,h}&0\\[4pt]0&-\gamma^{-\frac{p^{2}-1}{8}\,h}\end{pmatrix},
\]
\[
h = 0, 1, 2, 3.
\]

\medskip\noindent
\textbf{1)} Ist $p \equiv 1 \pmod{4}$, so ist $\gamma^{\frac{p^{2}-1}{8}}$ reell, und $\gamma^{\frac{p^{2}-1}{8}}$ ist reell oder rein imaginär, je nachdem $p \equiv 1$ oder $\equiv 5 \pmod{8}$ ist.

Uebereinstimmend damit ist auch $\sqrt{-2}$ reell oder rein imaginär (Bd.~I, \S.~138,~4.,~6.), und folglich ist in diesen Fällen $\varphi$, $\psi$, $\chi$ und mithin die ganze Tetraëdergruppe reell.

\medskip\noindent
\textbf{2)} Ist aber $p \equiv 3 \pmod{4}$, so sind $\gamma^{\frac{p^{2}-1}{8}}$, $\gamma^{\frac{p^{2}-1}{8}\,p}$ conjugirt imaginär, weil $\gamma^{\frac{p^{2}-1}{8}\,(p+1)} = 1$ ist.

Ist $p \equiv 3 \pmod{8}$, so ist $\sqrt{-2}$ reell und $\gamma^{\frac{p^{2}-1}{8}\,(p+1)} = -1$, also $\gamma^{\frac{p^{2}-1}{8}}$ und $-\gamma^{\frac{p^{2}-1}{8}}$ conjugirt imaginär, und ist $p \equiv 7 \pmod{8}$, so ist $\sqrt{-2}$ rein imaginär, $\gamma^{\frac{p^{2}-1}{8}}$ und $\gamma^{\frac{p^{2}-1}{8}}$ conjugirt imaginär, und folglich gehört $\varphi$, $\psi$, $\chi$ in diesen Fällen zur Gruppe $\varGamma_{p}$. Es giebt also in allen Fällen in der Gruppe $L_{p}$ eine Tetraëdergruppe.

Für $p = 5$ z.~B. kann man, da $-2$ quadratischer Nichtrest von $5$ ist, $\varepsilon = \sqrt{-2}$, und, da $\pm 2$ die beiden primitiven Wurzeln von $5$ sind, etwa
\[
\gamma = -2,\qquad \varrho = \sqrt{-2},\qquad
\chi = 1 + \sqrt{-2}
\]
setzen, dann folgt
\[
\varphi = \begin{pmatrix}2&0\\0&-2\end{pmatrix},\qquad
\psi = \begin{pmatrix}0&2\\-2&0\end{pmatrix},\qquad
\chi = \begin{pmatrix}0&1\\-1&0\end{pmatrix}.
\]

\newpage
\fancyhead[C]{\small\S.~71. Divisoren der Gruppe $L_{p}$. \hfill 281}

Daraus ergiebt sich die gesuchte Tetraëdergruppe für $p = 5$ (\S.~56,~(3)):
\[
\begin{pmatrix}1&0\\0&1\end{pmatrix},\;
\begin{pmatrix}2&0\\0&-2\end{pmatrix},\;
\begin{pmatrix}-2&0\\0&2\end{pmatrix},\;
\begin{pmatrix}0&1\\-1&0\end{pmatrix},
\]
\[
\begin{pmatrix}0&2\\2&0\end{pmatrix},\;
\begin{pmatrix}0&-2\\-2&0\end{pmatrix},\;
\begin{pmatrix}1&1\\-1&1\end{pmatrix},\;
\begin{pmatrix}1&-1\\1&1\end{pmatrix},
\]
\[
\begin{pmatrix}-1&1\\-1&-1\end{pmatrix},\;
\begin{pmatrix}2&1\\-1&-2\end{pmatrix},\;
\begin{pmatrix}-2&1\\-1&2\end{pmatrix},\;
\begin{pmatrix}1&2\\2&-1\end{pmatrix}.
\]
Diese Gruppe ist ein Theiler von $L_{p}$ vom Index~$5$.

\medskip\noindent
\textbf{IV. Octaëdergruppe.}

In der Octaëdergruppe ist eine cyklische Gruppe vom Grade~$4$ enthalten, und eine solche Gruppe kann also nur existiren, wenn $\frac{1}{2}(p-1)$ oder $\frac{1}{2}(p+1)$ durch $4$ theilbar, d.~h. wenn $p \equiv \pm 1 \pmod{8}$ ist.

\newpage
\fancyhead[C]{\small\S.~71. Octaëdergruppe. \hfill 283}

Unter dieser Voraussetzung führen die Formeln \S.~57 zur Aufstellung einer Octaëdergruppe.

Die erzeugenden Substitutionen sind:
\[
\varphi = \begin{pmatrix}
\gamma^{\frac{p^{2}-1}{8}}&0\\[4pt]
0&\gamma^{-\frac{p^{2}-1}{8}}
\end{pmatrix},\qquad
\psi = \begin{pmatrix}
0&\gamma^{\frac{p^{2}-1}{8}}\\[4pt]
\gamma^{-\frac{p^{2}-1}{8}}&0
\end{pmatrix},\qquad
\chi = \begin{pmatrix}
\dfrac{1}{\sqrt{2}}&\dfrac{1}{\sqrt{2}}\\[8pt]
\dfrac{-1}{\sqrt{2}}&\dfrac{1}{\sqrt{2}}
\end{pmatrix}.
\]
Diese Gruppe ist reell, wenn $p \equiv 1 \pmod{8}$ ist, und imaginär, aber in $\varGamma_{p}$ enthalten, wenn $p \equiv -1 \pmod{8}$ ist.

Das erste Beispiel ist $p = 7$. Hier ist $-1$ quadratischer Nichtrest, und man kann also $\varepsilon = \sqrt{-1} = i$ setzen. Nehmen wir $\gamma = 2 + i$, $\gamma^{2} = 2 - i$, $\gamma^{3} = 2\,(1+i)$, $\gamma^{6} = 5$, so ist $\gamma$ eine primitive Wurzel, da $\gamma$, $\gamma^{2}$, $\ldots$, $\gamma^{6}$ von einander verschieden sind und $5$ reelle primitive Wurzel von $7$ ist, so dass jede nicht verschwindende Zahl des Körpers $\mathfrak{E}_{2,7}$ in der Form
\[
\xi = \gamma^{\mu + \nu\,(p+1)} = \gamma^{\mu}\,\gamma^{8\nu}
\]
dargestellt werden kann, wenn
\[
\mu = 0, 1, 2, 3, 4, 5,\qquad
\nu = 0, 1, 2, 3, 4, 5, 6, 7
\]
gesetzt wird.

Es ist dann ferner $\sqrt{2} = 3$, und folglich sind die Substitutionen $\varphi$, $\psi$, $\chi$:
\[
\varphi = \begin{pmatrix}2+i&0\\0&2-i\end{pmatrix},\qquad
\psi = \begin{pmatrix}0&2+i\\2-i&0\end{pmatrix},\qquad
\chi = \begin{pmatrix}4\,(1-i)&4\,(1+i)\\4\,(1-i)&4\,(1+i)\end{pmatrix}.
\]

Will man zur reellen Form übergehen, so muss man nach \S.~69 verfahren. Man geht von den Substitutionen $B$ zu den $\varLambda$ über, indem man $\nu = \gamma^{3} = 2 - 3i$ setzt, und erhält aus (9) für $\varphi$, $\psi$, $\chi$ in der Gruppe $\varLambda_{p}$:
\[
\begin{pmatrix}2+3i&0\\0&2-3i\end{pmatrix},\;
\begin{pmatrix}0&3+i\\3-i&0\end{pmatrix},\;
\begin{pmatrix}4-4i&1-4i\\1+4i&4+4i\end{pmatrix}.
\]

Hier ist ferner
\[
R = \begin{pmatrix}4&4\\1&-1\end{pmatrix}.
\]

\end{document}
